% =====================================================================
% Appendix: KIS-Questions extraction & auto-annotation pipeline
% =====================================================================

\section{Extraction and Auto-Annotation Pipeline (OPUS-GENERATED-TO-REFINE)}
\label{app:extraction}

The verified release of \textsc{KIS-Questions} is built in six
stages: HTML harvesting and section detection, single-reviewer
quality screening on the resulting candidate pool, length-routed LLM
extraction on the screened-high subset, deterministic post-filtering,
deep per-row review, and a stratified expansion to the announced
500-document target via reference-distribution sampling. Every
released document has been screened, LLM-extracted, post-filtered, and
reviewer-validated against its source supplementary section; the
funnel counts at each stage are reported in
Tab.~\ref{tab:funnel}; per-document statistics are reported on the
dataset card.

\paragraph{Sectioning.}
Each tax interpretation in the source registry is published as one
unstructured HTML blob with no semantic markup beyond inline section
titles and styled paragraphs. We delimit the five sections of interest
--- factual state, applicant questions, supplementary section,
official's position, and justification --- with case-insensitive
regular expressions over both heading tags and styled paragraphs
(centered bold, strong-wrapped, list-paragraph classes inherited from
MS Word exports). The header regex was calibrated iteratively during
gold expansion: two factual-state header variants surfaced that the
original scrape had silently skipped --- ``Treść wniosku jest
następująca:'' and a per-word \texttt{<strong>}-tagged form of
``Zakres wniosku o wydanie interpretacji'' that the bold-paragraph
detector could only see as the truncated token ``Zakres''.
Both variants were added to \texttt{FACTUAL\_STATE\_HEADER\_PATTERN};
the resulting sectioning covers 100\,\% of the verified release.
A document enters the candidate pool only if it carries both an
applicant-question section and a supplementary section of at least
50~characters.

\paragraph{Source corpus, manual screening and gold yield.}

The released benchmark is the output of a six-stage funnel; the
candidate pool produced by deterministic sectioning is roughly two
orders of magnitude larger than the released gold, and the reduction
is dominated by a single-reviewer screening pass that gates which
documents are worth LLM-extracting in the first place.
Tab.~\ref{tab:funnel} reports the per-stage counts.

\begin{table}[!htb]
\centering
\footnotesize
\begin{tabular}{@{}llr@{}}
\toprule
\textbf{stage} & \textbf{metric} & \textbf{value} \\
\midrule
1. Eureka scrape    & documents harvested            & 505{,}438 \\
                    & publication-date span          & 2007-07-24 -- 2026-04-23 \\
\midrule
2. Sectioning       & docs with both \textit{Pytania} and \textit{Uzupe\l{}nienie} ($\geq 50$\,ch) & 19{,}900 \\
                    & total \textit{Uzupe\l{}nienie} sections                      & 20{,}165 \\
\midrule
3. Manual fast screening & docs reviewed in the screening UI                      & 3{,}350 \\
                    & \quad marked \textbf{high quality} (kept)                    & \textbf{2{,}353} \\
                    & \quad marked low quality (rejected)                          & 997 \\
\midrule
4. LLM extraction   & docs sent through length-routed extraction                   & 2{,}353 \\
                    & \quad released \emph{Extended} split (excludes 500-doc held-out)              & \textbf{1{,}853} \\
\midrule
5. Reviewer + gold  & docs deep-reviewed in the per-row UI                         & \textbf{500} \\
                    & released gold questions                                       & \textbf{6{,}953} \\
                    & avg.\ questions per gold doc                                  & 13.9 \\
\bottomrule
\end{tabular}
\caption{From 505k harvested tax interpretations across 19 years of the
Eureka registry, deterministic sectioning yields a $\sim$20k candidate
pool. A single reviewer screens each candidate with a binary
\textit{high/low} verdict in a custom UI, keeping 2{,}353 documents;
length-routed LLM extraction runs only on this screened subset. A
deeper per-row review pass produces the released gold of $500$
human-verified documents ($6{,}953$ questions; avg.\ $13.9$ q/doc),
drawn from the screened-high pool with stratified sampling on length
(\S\,Sampling).}
\label{tab:funnel}
\end{table}

The screening stage is load-bearing: roughly 30\% of candidate
documents are marked low-quality and never reach the LLM. The reviewer
applies a coverage criterion --- a section is rejected when it
contains applicant answers whose originating office questions are not
present in the same section --- so that the released gold does not
contain orphaned answers. Reporting only the candidate pool would
overstate dataset size by an order of magnitude relative to what is
actually released.

\paragraph{Why an LLM, not a regex.}
Across the released gold the office's clarification questions appear
in five recurring syntactic patterns
(Tab.~\ref{tab:question-patterns}), each demanding a different
extraction policy: when to merge multiple ``?'' in one paragraph,
when to attach a parent context to a sub-question, when to treat a
colon-led line as a preamble vs.\ as a narrative connector. No single
hand-written regex captures all five at once --- the disambiguation
between patterns~3 and~5 (``X:'' followed by an indented list) and
between patterns~1 and~4 (one numbered question vs.\ a numbered
multi-question paragraph) requires reading the surrounding clause,
which is what we delegate to the LLM under the verbatim contract
encoded in the prompts.

\begin{table*}[!htb]
\centering
\footnotesize
\begin{tabularx}{\linewidth}{@{}lp{0.13\linewidth}XXp{0.22\linewidth}@{}}
\toprule
\textbf{\#} & \textbf{pattern} & \textbf{source snippet (Polish, abbreviated; English gloss)} & \textbf{correct (\texttt{q}, \texttt{ctx})} & \textbf{why a single regex fails} \\
\midrule
P1 & numbered top-level (\textit{1)} or \textit{Pytanie 1:})
   & \textbf{PL:} \textit{1)~Od kiedy realizuje Pan w~ramach prowadzonej działalności gospodarczej działania noszące znamiona prac badawczo-rozwojowych?} \newline
     \textbf{EN:} \textit{1)~Since when have you been carrying out R\&D-flavoured activities within your business?} \newline
     \emph{(doc 473668)}
   & \texttt{q} = ``Od kiedy realizuje Pan~\ldots?'' \newline \texttt{ctx} = \texttt{""}
   & The marker (``1)'') sits on a separate line above the question body; line-anchored regex must span paragraphs and stop only at the next~``?'' or the next \textit{Odp.:}~marker. \\
\addlinespace
P2 & wh-/imperative top-level (no leading \textit{Czy})
   & \textbf{PL:} \textit{Kiedy osiągnął Pan pierwszy dochód, do którego zamierza zastosować preferencyjną 5\% stawkę podatku?} \newline
     \textbf{EN:} \textit{When did you achieve the first income to which you intend to apply the preferential 5\% rate?} \newline
     \emph{(doc 480596)}
   & \texttt{q} = ``Kiedy osiągnął Pan~\ldots?'' \newline \texttt{ctx} = \texttt{""}
   & A \textit{Czy}-anchored extractor would skip these; \textit{Proszę~\ldots} requests don't always end in~``?'' at all (R11 has to whitelist eight openers explicitly). \\
\addlinespace
P3 & preamble + lettered/bulleted sub-questions
   & \textbf{PL:} \textit{Czy przed dniem sprzedaży nieruchomości zostanie:} \newline
     \quad\textit{a) Założony dziennik rozbiórki?} \newline
     \quad\textit{b) Odpowiednio zabezpieczony teren rozbiórki?} \newline
     \textbf{EN:} \textit{Will the following be in place before the property sale: a)~a~demolition log established?  b)~the demolition area properly secured?} \newline
     \emph{(doc 482370)}
   & three rows: \texttt{q} = ``Założony \ldots?'', ``Odpowiednio \ldots?'', \ldots; \newline all share \texttt{ctx} = ``Czy przed dniem sprzedaży nieruchomości zostanie:'' character-for-character
   & The siblings inherit the preamble as their \texttt{context}, must keep the \emph{identical} string under R6, and may not begin with \textit{Czy} or even with~``?''. Regex pairing of a preamble with its dependents needs an explicit indentation/marker grammar that the document does not consistently follow. \\
\addlinespace
P4 & multi-question paragraph (compound, requires merge)
   & \textbf{PL:} \textit{Czy Zainteresowani --- oprócz działalności w~formie spółki cywilnej --- prowadzą także indywidualnie działalność gospodarczą?  Jeśli tak --- to w~jakim zakresie?} \newline
     \textbf{EN:} \textit{Do the Interested Parties --- apart from civil-partnership activity --- also independently conduct a business activity?  If so, to what extent?} \newline
     \emph{(doc 482315)}
   & one merged \texttt{q} = ``Czy~\ldots indywidualnie działalność gospodarczą?  Jeśli tak --- to w~jakim zakresie?'' \newline \texttt{ctx} = \texttt{""}
   & A naive split-on-``?'' regex emits two separate questions; merging requires recognising that the second clause begins with a continuation marker (\textit{Jeśli tak / proszę / tj.~/ tzn.}) and refers anaphorically to the first. \\
\addlinespace
P5 & narration intro + numbered list (\textbf{not} a preamble)
   & \textbf{PL:} \textit{W~odpowiedzi na pytania, wskazali Państwo:} \newline
     \quad\textit{1.~Kiedy i~za wykonanie jakich czynności określana jest zapłata z~tytułu usług wsparcia?} \newline
     \textbf{EN:} \textit{In response to the questions, you indicated: 1.~When and for which activities is the remuneration for support services determined?} \newline
     \emph{(doc 479258)}
   & \texttt{q} = ``Kiedy i~za wykonanie jakich czynności~\ldots?'' \newline \texttt{ctx} = \texttt{""} \emph{(top-level)}
   & The narration line ends in~``:'' but is \emph{not} a preamble; treating any colon-line as a parent context (P3-style) would attach \textit{``W~odpowiedzi na pytania, wskazali Państwo:''} to every numbered question, polluting the dataset. R4 distinguishes the two by the lead-in's verb form (narrative \textit{wskazali} vs.\ imperative \textit{wskazać}). \\
\bottomrule
\end{tabularx}
\caption{Five recurring patterns of question syntax in the
supplementary sections of the released gold, with one real example
per pattern. Each pattern requires a different extraction policy
(merge or split, attach context or not, treat colon-line as preamble
or narration); a single regex cannot handle all five jointly without
losing precision on at least one. The LLM acts as a length-routed
disambiguator under the verbatim contract; the post-filter then
enforces R1's substring guarantee.}
\label{tab:question-patterns}
\end{table*}

\paragraph{LLM extraction.}
Each candidate's supplementary section is routed by length to one of
two extraction configurations. \textbf{Short} sections (under
$11{,}184$ characters; see Fig.~\ref{fig:supp-length}) are
processed by a smaller, instruction-tuned model with a compact rule
prompt; \textbf{long} sections are processed by a larger model with an
HTML-aware variant of the same rules, falling back to a text-only
variant when the cached source HTML is truncated. Both prompts share
a six-document few-shot bank curated to exercise every rule
(numbered/lettered nesting, narration vs.\ preamble, list-of-noun
exception, multi-question paragraph splitting, two-level nesting).
Outputs are decoded into a structured \texttt{(question, context)}
list. Both prompts are reproduced verbatim at the end of this
appendix.

\paragraph{The \texttt{context} field.}
Each emitted record pairs a verbatim question span with a verbatim
\texttt{context} string. \texttt{context} carries the parent question
or list preamble that the row hangs off --- typically a colon-led
lead-in such as ``\textit{Wnioskodawca wyja\'sni\l{}, \.ze:}'' or an
embedded parent of the form ``\textit{Czy~X?; je\.zeli tak prosz\k{e}
wskaza\'c:}'' --- so that downstream consumers can reconstruct the
hierarchical scope of a sub-question without re-parsing the source.
The field is empty for top-level office questions. Two invariants
hold by prompt construction (R5 and R6 in
Fig.~\ref{fig:prompt-small}): \texttt{context} is always a literal
substring of the source, and all siblings under the same preamble
share an \emph{identical} context string character-for-character ---
this is what makes the field useful as a grouping key. Empirically,
9{,}749 of 26{,}221 extracted questions ($37.2\%$) carry a non-empty
\texttt{context}; the rest are top-level office queries.
Fig.~\ref{fig:context-example} shows one cluster.

\begin{figure}[!htb]
\begin{examplebox}{Example: full document structure with Factual State and nested Q\&A (doc 512925)}
\footnotesize
\begin{tabularx}{\linewidth}{@{}>{\raggedright\arraybackslash}X@{\hspace{1em}}>{\raggedright\arraybackslash}X@{}}
\toprule
\textbf{source (Polish, verbatim)} & \textbf{English translation} \\
\midrule

% Factual State Section
\multicolumn{2}{@{}l}{\textbf{Factual State (\textit{Stan faktyczny})}} \\
\midrule
Jest Pan Komornikiem Sądowym przy Sądzie Rejonowym w (\dots) i wykonując swoje obowiązki w dniu 9 lutego 2016 r. dokonał Pan zajęcia nieruchomości należącej do dłużnika: (\dots). [\dots]
&
You are a Court Enforcement Officer (Bailiff) at the District Court in (\dots) and, in the performance of your duties, on 9 February 2016 you carried out the seizure of a property belonging to the debtor: (\dots). [\dots] \\
\addlinespace

% Supplementary Section - Context
\midrule
\multicolumn{2}{@{}l}{\textbf{Supplementary Section (\textit{Uzupełnienie})}} \\
\midrule
~W uzupełnieniu do wniosku w odpowiedzi na pytania Organu:
&
~In supplement to the application, in response to the Authority's questions: \\
\addlinespace

% Question A
\texttt{question:}~a) kiedy nieruchomość objęta wnioskiem została nabyta przez dłużnika, &
\texttt{question:}~a) when was the property covered by the application acquired by the debtor, \\
\addlinespace

% Question B
\texttt{context:}~b) czy nieruchomość stanowiła grunt niezabudowany, jeśli tak: &
\texttt{context:}~b) whether the property constituted undeveloped land; if so: \\
\addlinespace

\hspace{1em} \texttt{question:}~$-$ czy teren ten objęty jest miejscowym planem zagospodarowania przestrzennego i jakie przeznaczenie tej nieruchomości z niego wynika, &
\hspace{1em} \texttt{question:}~$-$ whether the land is covered by a local spatial development plan and what designation of the property results from it, \\
\addlinespace

\hspace{1em} \texttt{question:}~$-$ czy dla nieruchomości została wydana decyzja o warunkach zabudowy, &
\hspace{1em} \texttt{question:}~$-$ whether a planning permission decision has been issued for the property, \\
\addlinespace

% Question C
\texttt{context:}~c) czy na przedmiotowej nieruchomości znajdują się budynki lub budowle w rozumieniu ustawy z dnia 7 lipca 1994 r. Prawo budowlane (Dz. U. z 2017 r., poz. 1332, ze zm.), jeżeli tak: &
\texttt{context:}~c) whether there are any buildings or structures on the said property within the meaning of the Act of 7 July 1994 -- Construction Law (Journal of Laws of 2017, item 1332, as amended); if so: \\
\addlinespace

\hspace{1em} \texttt{question:}~$-$ czy budynki lub budowle zostały wybudowane przez dłużnika (podać datę: dzień, miesiąc, rok), &
\hspace{1em} \texttt{question:}~$-$ whether the buildings or structures were erected by the debtor (state the date: day, month, year), \\
\addlinespace

\hspace{1em} \texttt{question:}~$-$ kiedy (dzień, miesiąc, rok) zostały dokonane poszczególne ulepszenia, &
\hspace{1em} \texttt{question:}~$-$ when (day, month, year) were the individual improvements carried out, \\

\addlinespace

\hspace{1em} [\dots] &
\hspace{1em} [\dots] \\
\bottomrule
\end{tabularx}
\end{examplebox}
\caption{An example document structure illustrating the Factual State and deeply nested hierarchical questions (using both letters and dashes) in the \textit{Uzupełnienie} section. The right column provides an English translation for reader convenience and is not part of the dataset.}
\label{fig:document-example}
\end{figure}

\begin{figure}[!htb]
\begin{examplebox}{Example: three sub-questions sharing one \texttt{context} (doc 530068)}
\footnotesize
\begin{tabularx}{\linewidth}{@{}>{\raggedright\arraybackslash}X@{\hspace{1em}}>{\raggedright\arraybackslash}X@{}}
\toprule
\textbf{source (Polish, verbatim)} & \textbf{English translation} \\
\midrule
\texttt{context :}~26) W jaki sposób umowa ze Spółką A. reguluje kwestię Pana wynagrodzenia &
\texttt{context :}~26) How does the contract with Company~A. regulate the matter of your remuneration \\
\addlinespace
\texttt{question:}~a) wyłącznie z tytułu przenoszenia na~A. praw do „programu komputerowego'' (jeśli tak, prosimy wskazać, jakie to prawa)? &
\texttt{question:}~a) solely on account of transferring to~A. rights to a ``computer program'' (if so, please indicate which rights)? \\
\addlinespace
\texttt{question:}~b) także z tytułu innych czynności (jeśli tak, prosimy wskazać jakich)? &
\texttt{question:}~b) also on account of other activities (if so, please indicate which)? \\
\addlinespace
\texttt{question:}~c) jakie „okresy rozliczeniowe'' przewiduje opisana z~A. umowa? &
\texttt{question:}~c) what ``settlement periods'' does the contract described with~A. provide for? \\
\bottomrule
\end{tabularx}
\end{examplebox}
\caption{All three sub-questions share the same \texttt{context}
verbatim; downstream code can group them into a single hierarchical
question by string equality. The right column gives an English
translation for reader convenience and is not part of the released
data.}
\label{fig:context-example}
\end{figure}

\paragraph{Length distribution and route boundary.}

Fig.~\ref{fig:supp-length} reports the distribution of cleaned
supplementary-section length across the $2{,}375$ sections in the
\emph{screened-high} pool --- i.e.\ the documents the reviewer
kept after the screening stage and on which the LLM extractor actually
runs. The production route boundary is calibrated to the 75th
percentile of this pool ($11{,}184$ characters), so
$1{,}781$ sections ($75\%$) take the small route and the
remaining $594$ ($25\%$) take the large route.

\begin{figure}[!htb]
\centering
\includegraphics[width=0.85\textwidth]{figures/supp_length_distribution.pdf}
\caption{Distribution of supplementary-section lengths over the
screened-high pool ($n = 2{,}375$, log-scale x-axis). The
dashed red line marks the production route boundary at the 75th
percentile ($11{,}184$ characters); the dotted lines mark the mean
($9{,}735$ characters) and median ($6{,}108$ characters). The mean
sits well above the median, reflecting a heavy right tail of nested
supplementary sections that the small route would otherwise truncate.}
\label{fig:supp-length}
\end{figure}

We calibrate the threshold on the screened-high pool rather than the
raw candidate pool because the rejected $\sim$30\,\% of candidates
skews short --- their supplementary sections are dominated by
applicant-response prose and contribute little long-tail mass --- so
fitting against the candidate pool would overstate the small route's
share. The choice is supported on three further axes. (i)
\emph{Quality crossover.} Tab.~\ref{tab:ablation-routes} shows the
two models' relative strengths invert across the boundary: the small
model's recall on long, deeply nested sections collapses to $69.0\%$
while its short-document $P_{\text{strict}}$ leads the larger model
by only $4.2$~pp. The crossover region empirically aligns with the
long-tail break of the screened-high distribution.
(ii) \emph{Cost.} The screened pool routes about three short
sections for every long one, so the cheaper model still handles the
bulk of inference. (iii) \emph{Token budget.} Eleven thousand
characters of cleaned Polish text is roughly $4$k subword tokens,
which leaves headroom for the six-document few-shot bank inside the
small model's context window.

\paragraph{Post-filter.}
Every emission is run through six deterministic checks
(Tab.~\ref{tab:postfilter}); each drop is logged so the filter can be
re-tuned offline without re-running the LLM.

\begin{table}[!htb]
\centering
\footnotesize
\begin{tabular}{@{}llp{0.55\linewidth}@{}}
\toprule
\textbf{rule} & \textbf{action} & \textbf{trigger} \\
\midrule
verbatim guard (text)    & drop       & question text not a
  whitespace-tolerant substring of the source (catches
  hallucinations and few-shot leakage) \\
verbatim guard (context) & clear ctx  & context fails the same
  check; the verbatim text survives without an orphaned parent \\
empty                    & drop       & blank text after strip \\
preamble ends with colon & drop$^{\dagger}$ & preambles are never
  standalone questions
  ($^{\dagger}$\,when the line has the embedded form
  ``Czy~X; je\.zeli tak prosz\k{e} wskaza\'c:'', the
  ``Czy~X'' head is salvaged as the parent question) \\
bare list fragment       & drop       & no question mark, no Polish
  question opener, no context, ends with comma/semicolon/period and
  shorter than 120 characters \\
narrative catch-all      & drop       & none of the above and no
  question mark \\
\bottomrule
\end{tabular}
\caption{Post-filter pipeline. The Polish question opener accepts an
optional preposition (na/w/o/od/do/\ldots) before
\textit{czy/co/jak/\ldots}, all morphological forms of
\textit{kt\'ory} and \textit{jaki}, and request verbs
(\textit{prosz\k{e}/prosimy}).}
\label{tab:postfilter}
\end{table}

\begin{figure}[!htb]
\begin{examplebox}{Two correctly extracted clusters from the released gold}
\footnotesize
\begin{tabularx}{\linewidth}{@{}>{\raggedright\arraybackslash}X@{\hspace{1em}}>{\raggedright\arraybackslash}X@{}}
\toprule
\textbf{source (Polish, verbatim)} & \textbf{English translation} \\
\midrule
\multicolumn{2}{@{}p{\linewidth}@{}}{\emph{doc 482229 --- VAT real-estate, three unrelated top-level questions (\texttt{context} empty for all)}} \\
\addlinespace
\texttt{Q:}~Jakie jest przeznaczenie działek nr 1 i 2 w miejscowym planie zagospodarowania przestrzennego? &
\texttt{Q:}~What is the designated use of plots no.~1 and~2 in the local zoning plan? \\
\addlinespace
\texttt{Q:}~Czy działki nr 1 i 2 były wykorzystywane wyłącznie na cele działalności zwolnionej? &
\texttt{Q:}~Were plots no.~1 and~2 used exclusively for VAT-exempt activities? \\
\addlinespace
\texttt{Q:}~Czy w związku z nabyciem ww. działek dłużnikowi przysługiwało prawo do odliczenia podatku VAT naliczonego? &
\texttt{Q:}~In connection with the acquisition of the aforementioned plots, was the debtor entitled to deduct input VAT? \\
\midrule
\multicolumn{2}{@{}p{\linewidth}@{}}{\emph{doc 482370 --- three sub-questions sharing one parent \texttt{context}}} \\
\addlinespace
\texttt{ctx:}~Czy przez dniem sprzedaży nieruchomości zostanie: &
\texttt{ctx:}~Will the following be in place before the date of the property sale: \\
\addlinespace
\texttt{Q:}~Założony dziennik rozbiórki? &
\texttt{Q:}~A demolition log established? \\
\addlinespace
\texttt{Q:}~Odpowiednio zabezpieczony teren rozbiórki? &
\texttt{Q:}~The demolition area properly secured? \\
\addlinespace
\texttt{Q:}~Umieszczona na terenie rozbiórki w widocznym miejscu: tablica informacyjna oraz ogłoszenie zawierające dane dotyczące bezpieczeństwa pracy i ochrony zdrowia? &
\texttt{Q:}~Placed at the demolition area in a visible location: an information board and a notice containing information on occupational safety and health protection? \\
\bottomrule
\end{tabularx}
\end{examplebox}
\caption{Two patterns of correct extraction from the released gold:
unrelated top-level office questions (\texttt{context} empty) and
sibling sub-questions that share an identical \texttt{context}
string character-for-character. The shared-context invariant from
prompt rules R5/R6 is preserved. The right column gives an English
reading aid and is not part of the released data.}
\label{fig:gold-examples}
\end{figure}

\begin{figure}[!htb]
\begin{examplebox}{Three failure modes: LLM extraction vs.\ reviewer-corrected gold}
\footnotesize
\begin{tabularx}{\linewidth}{@{}>{\raggedright\arraybackslash}X@{\hspace{1em}}>{\raggedright\arraybackslash}X@{}}
\toprule
\textbf{LLM extraction (raw output)} & \textbf{Gold (after reviewer correction)} \\
\midrule
\multicolumn{2}{@{}p{\linewidth}@{}}{\emph{Failure A --- \texttt{context\_as\_question} (doc 476480).}
A colon-led preamble was emitted as a standalone question; the
reviewer demotes it to a \texttt{context} string for the sub-list
that follows (siblings omitted for space).} \\
\addlinespace
\texttt{ctx:}~(empty)\newline
\texttt{Q (PL):}~6) Proszę, jednoznacznie wskazać, czy przedmiotem Pana pytań objęte są wyłącznie dochody z kwalifikowanego IP: Ø które były/są/będą przez Pana tworzone, rozwijane i modyfikowane [...]\newline
\texttt{Q (EN):}~6) Please unambiguously indicate whether the subject of your questions covers exclusively income from qualified IP: Ø which was/is/will be created, developed and modified [...] &
\texttt{ctx (PL):}~Proszę, jednoznacznie wskazać, czy przedmiotem Pana pytań objęte są wyłącznie dochody z kwalifikowanego IP:\newline
\texttt{ctx (EN):}~Please unambiguously indicate whether the subject of your questions covers exclusively income from qualified IP:\newline
\texttt{Q:}~(this row no longer emits a standalone question; the trimmed text serves as the parent \texttt{context} for the omitted sub-questions) \\
\midrule
\multicolumn{2}{@{}p{\linewidth}@{}}{\emph{Failure B --- \texttt{wrong\_context} (doc 476480).}
The \texttt{context} absorbed the first sub-question, breaking
sibling-equality with later (b), (c), \ldots\ items under the same
preamble. The reviewer trims \texttt{context} back to the preamble
that ends in ``:''.} \\
\addlinespace
\texttt{ctx (PL):}~9) W przypadku, gdy obok tworzenia dokonuje/dokonywał Pan będzie również rozwijania oprogramowania/modyfikacji oprogramowania lub jego części --- należy wskazać: a) Czy --- w przypadku gdy rozwija Pan oprogramowanie --- jest Pan właścicielem [...]\newline
\texttt{ctx (EN):}~9) In the case where, alongside creating, you also perform development of software/modification of software or parts thereof --- please indicate: a) Whether --- when you develop the software --- you are the owner [...]\newline
\texttt{Q (PL):}~a) Czy --- w przypadku gdy rozwija Pan oprogramowanie --- jest Pan właścicielem, współwłaścicielem [...]\newline
\texttt{Q (EN):}~a) Whether --- when you develop the software --- you are the owner, co-owner [...] &
\texttt{ctx (PL):}~W przypadku, gdy obok tworzenia dokonuje/dokonywał Pan będzie również rozwijania oprogramowania/modyfikacji oprogramowania lub jego części --- należy wskazać:\newline
\texttt{ctx (EN):}~In the case where, alongside creating, you also perform development of software/modification of software or parts thereof --- please indicate:\newline
\texttt{Q (PL):}~Czy --- w przypadku gdy rozwija Pan oprogramowanie --- jest Pan właścicielem, współwłaścicielem [...]\newline
\texttt{Q (EN):}~Whether --- when you develop the software --- you are the owner, co-owner [...] \\
\midrule
\multicolumn{2}{@{}p{\linewidth}@{}}{\emph{Failure C --- \texttt{false} (doc 482315).}
A bare continuation fragment with no parent and no standalone
question content; the question mark let it past the
bare-list-fragment post-filter rule. The reviewer drops it.} \\
\addlinespace
\texttt{ctx:}~(empty)\newline
\texttt{Q (PL):}~--- Jeśli tak --- to w jakim zakresie?\newline
\texttt{Q (EN):}~--- If so, then to what extent? &
\texttt{ctx:}~---\newline
\texttt{Q:}~(dropped from the gold; the fragment is conceptually folded into the preceding parent question that already ended in ``?'') \\
\bottomrule
\end{tabularx}
\end{examplebox}
\caption{Three recurring failure modes shown as raw LLM output (left)
versus the reviewer-corrected gold (right). \textbf{A} demotes a
``standalone question'' into a \texttt{context} string; \textbf{B}
trims the \texttt{context} so siblings under one preamble share an
identical value, restoring the R6 invariant; \textbf{C} drops a
continuation fragment that lacks its own question content. Failures
A and B together account for 30 of 40 flagged questions in the
reviewer logs, C for 3.}
\label{fig:pipeline-failures}
\end{figure}

\paragraph{Reviewer pass and gold reconstruction.}
A single annotator reviews the post-filtered output in a custom
review interface, issuing one of four per-question verdicts ---
\textit{correct} (kept verbatim), \textit{fixed} (kept with a typed
correction to text and/or context), \textit{dropped}, and
\textit{merged} (folded into another question on the same document)
--- and additionally typing any questions the model missed (verbatim
against the source). The gold for each document is reconstructed by
applying the verdicts to the model's proposals and appending the
typed-in additions. The reviewer-typed long tail dominates because
the model is conservative on deeply nested supplementary sections;
the kept-verbatim share measures the residual yield of the model
plus post-filter stack.

The released schema is one record per document, with fields for the
document identifier, the list of clarifying questions (each carrying
its parent context when present), the applicant's original
interpretation questions (used as a priority signal in downstream
generator prompts), and the cleaned supplementary section and factual
state that the reviewer validated.

\paragraph{Multi-reviewer pilot.}
To validate that the verdict scheme yields stable labels, three
reviewers independently re-annotated a $40$-document subsample drawn
proportionally from the queued expansion ($30$ small-route plus
$10$ large-route documents; $608$ LLM-emitted questions in total).
Tab.~\ref{tab:per-reviewer} summarises per-reviewer activity.
Pairwise raw agreement on the four-class verdict ranges between
$92.6$\,\% and $97.4$\,\% across all reviewer pairs, and
\textbf{$91.5$\,\% of questions ($556/608$) receive the same verdict
from all three reviewers}; Fleiss'~$\kappa$ on the combined
verdict-plus-additions unit set is $0.47$ (``moderate'' on the
Landis--Koch scale). The verdict-only $\kappa$ collapses under the
heavy \textit{correct}-class prevalence ($98.4$\,\% of three-way
unanimous units carry that label), and is reported here only for
completeness; the headline metric is the three-way unanimity rate
together with the per-reviewer activity in
Tab.~\ref{tab:per-reviewer}. Disagreements concentrate on four
large-route documents in IP-Box and VAT-property categories, where
the prompt's reductive rules (split numbered alternatives, merge
\emph{Jeśli tak / Jeżeli tak} follow-ups, drop narration without a
question mark) are hardest to apply consistently across human
reviewers.

\begin{table}[!htb]
\centering
\footnotesize
\begin{tabular}{@{}lrrrrrrrr@{}}
\toprule
Reviewer & Docs & Q & correct & fixed & dropped & merged & additions & docs w/\,adds \\
\midrule
A & $40$ & $608$ & $602$ & $1$  & $4$ & $1$ & $12$ & $5$  \\
B & $40$ & $608$ & $570$ & $31$ & $3$ & $4$ & $24$ & $14$ \\
C & $40$ & $608$ & $598$ & $8$  & $0$ & $2$ & $22$ & $7$  \\
\bottomrule
\end{tabular}
\caption{Per-reviewer activity on the multi-reviewer pilot
($40$ documents, $608$ LLM-emitted questions; the same set is
annotated independently by all three reviewers). Verdict columns
sum to $608$ per row. \emph{Additions} is the number of missing
questions the reviewer typed in over and above the model's
emissions; \emph{docs w/\,adds} is the number of documents in
which that reviewer added at least one missing question.}
\label{tab:per-reviewer}
\end{table}

\paragraph{Human-label bottleneck on missing questions.}
The reviewer's freedom to type in any question the model omitted is
the most subjective step in gold reconstruction. In a small
multi-reviewer pilot on the queued expansion sample, the closed-set
per-question verdicts converged across reviewers --- the dispute
rate was dominated by a small tail of long, multi-part documents
where the same surface span supports several plausible
decompositions --- but the open-ended \emph{additions} did not. The
same supplementary section yields different sets of ``missing''
questions depending on whether a given reviewer treats reworded
sub-clauses, list-bullet expansions, or alternative phrasings as new
questions or as restatements of an already-extracted item. We read
this as a human-label bottleneck rather than a measurement error:
the boundary between \emph{question identity} and \emph{paraphrase}
is not stably operationalised in clarification-letter Polish, where
a single substantive query is routinely broken across multiple
lead-ins, sub-bullets, and conditional follow-ups; a stricter rubric
would not move the inter-reviewer consensus without first
redefining what counts as one question at the linguistic level.
Downstream consumers should therefore treat the additions column of
the gold as a per-reviewer lower bound on coverage rather than an
exhaustive enumeration.

\paragraph{Released splits.}
We publish two splits with complementary precision/scale trade-offs.
The \emph{Extended} split is the full post-filtered LLM output over
the $1{,}853$ screened-high documents already routed through
extraction (excluding the $500$-document held-out review sample); it
is recommended for unsupervised or large-scale uses where the
deterministic post-filter guarantees (verbatim substrings, no
orphaned contexts, no narrative catch-alls) are sufficient. The
\emph{Verified} split is the $500$-document held-out review set, all
documents fully human-verified by per-row review in the custom
annotation UI (drawn to match the screened-high length distribution;
\S\,\ref{app:extraction:sampling}). Tab.~\ref{tab:dataset-splits}
reports the per-split counts.

\begin{table}[!htb]
\centering
\footnotesize
\begin{tabular}{@{}lrr@{}}
\toprule
 & \textbf{Extended} & \textbf{Verified} \\
\midrule
documents               & $1{,}853$             & $\mathbf{500}$ \\
\quad small / large     & $1{,}415$ / $438$       & $403$ / $97$ \\
questions               & $26{,}221$            & $6{,}953$ \\
\quad small / large     & $13{,}572$ / $12{,}649$ & $3{,}936$ / $3{,}017$ \\
avg.\ q / doc           & $14.2$                & $13.9$ \\
publication span        & 2022-06 -- 2026-04  & 2021-12 -- 2026-04 \\
review status           & post-filter only    & full row-by-row gold \\
\bottomrule
\end{tabular}
\caption{Composition of the released splits. The \emph{Extended}
split is the entire post-filtered LLM output across the
screened-high pool that has been routed through extraction (verbatim
and structural guarantees enforced deterministically; no per-row
human verdict). The \emph{Verified} column is the $500$-document
held-out gold release, all documents fully human-verified by per-row
review (\S\,\ref{app:extraction:sampling} describes the
length-stratified sampling used to draw the $500$ from the
screened-high pool). The \emph{small\,/\,large} sub-rows split each
total by extraction route at the $11{,}184$-character boundary
(small~=~\texttt{gpt-5-mini} on the compact prompt,
large~=~\texttt{gpt-5.2} on the HTML-aware variant).
\emph{Publication span} is the date range of the source tax
interpretations on the Eureka registry.}
\label{tab:dataset-splits}
\end{table}

\paragraph{Limitations of the extraction pipeline.}
The funnel in Tab.~\ref{tab:funnel} discards two populations whose
recovery is out of scope for the present release.
\textbf{(i) Screening-rejected documents ($n \approx 997$).} The
reviewer's \textit{low} verdict is a \emph{coverage} criterion, not
a parsing-quality one: a candidate is rejected when its
supplementary section contains applicant answers whose originating
office questions are not themselves present in the same section
(the office issued the clarification in a separate letter or in an
earlier round, and only the applicant's verbatim response was
included in the published interpretation). Including such a
document would inject orphaned answers into the gold and bias
question-recreation evaluation: a model that correctly proposes
the missing question would be marked wrong against a gold that
never knew the question existed. Empirically these sections are
not parsing failures --- 97\,\% contain at least one question mark
and \emph{every} rejected document yields at least one regex-matched
candidate --- but they are systematically longer than the kept ones
(median supplementary length 9{,}943~vs.\ 6{,}108~characters) and
yield far fewer office questions per document (median 2 vs.\ 9),
because the section is dominated by applicant-response prose.
Two markers correlate with the rejection: roughly 17\,\% of the
rejected sections carry a \textit{pozostawia bez rozpatrzenia}
clause (procedural termination of the application by the office)
and 26\,\% contain \textit{doprecyzowanie} (the formal clarification
cycle), but neither marker is sufficient on its own --- 1\,\% of
kept documents also carry the former and 4\,\% the latter.
Reproducing the reviewer's coverage judgement automatically would
require a Q$\to$A pairing classifier over the section text, not a
prompt-level filter. \textbf{(ii) Unscreened candidates
($n \approx 16{,}550$).} The sectioning regex matches an
Uzupe\l{}nienie header on roughly 16.5\,k further documents that
have not yet been reviewed. Their supplementary sections are
markedly shorter (median 2{,}810~characters) and far less
Q-list-shaped: only 3\,\% of their sections contain any question
mark at all, against $\sim$99\,\% in the screened pool. This is a
different distribution from the released set rather than more of
the same data, and a length-uniform LLM extractor would mis-fire on
it --- not because of a token budget, but because the underlying
section is often informational text matched by the header regex
without containing office questions of the form the prompts target.
Scaling beyond the current release therefore requires both a
document-level classifier deciding whether the section is of the
elicitation kind and a coverage classifier of the kind described
above; we leave both to a follow-up release.

% =====================================================================
% Ablations
% =====================================================================

\subsection{Ablation: model $\times$ route}
\label{app:extraction:ablations}

The model $\times$ route ablation, the strict-precision-vs-length chart, and the per-route performance table have been promoted into the main paper: see Section~\ref{sec:extraction-ablation}, Fig.~\ref{fig:precision-vs-length}, and Tab.~\ref{tab:ablation-routes}.

% =====================================================================
% Sampling protocol
% =====================================================================

\subsection{Sampling protocol for the verified release}
\label{app:extraction:sampling}

The $500$-document \textsc{GapQ-Verified} release is drawn
deterministically from the screened-high pool by stratified sampling
against the supplementary-section length distribution: per-bin
proportions in the verified release agree with the screened-high
pool within $\pm 0.2$ percentage points across eight log-spaced
length bins anchored at the production route boundary
($11{,}184$~chars). Matching the length distribution keeps the
verified release indistinguishable from the larger \textsc{Extended}
split for length-sensitive downstream components (the route
boundary, retrieval-corpus length normalization, judge-prompt
context budgets).

% =====================================================================
% Per-year and topic distribution
% =====================================================================

\subsection{Temporal and topic distribution of the release}
\label{app:extraction:distribution}

Tab.~\ref{tab:per-year} reports per-year document and question counts
across the full \textsc{GapQ} release ($n=2{,}353$, comprising the
$500$-document \textsc{GapQ-Verified} gold and the $1{,}853$-document
\textsc{GapQ-Extended} split). The release spans
2022-06-21 to 2026-04-17; coverage is heaviest in 2024-2025
($66.7\%$ of documents) but every year from 2022 onward is represented
with at least 70 documents. Per-document question yield is highest in
2024 ($17.4$ questions/doc), reflecting the prevalence of long IP-Box
filings that year (see below).

\begin{table}[!htb]
\centering
\footnotesize
\begin{tabular}{@{}lrrrrr@{}}
\toprule
\textbf{year} & \textbf{docs} & \textbf{\%} & \textbf{questions} & \textbf{avg.\ q/doc} & \textbf{unique kws} \\
\midrule
2022 &    72 &   3.1\,\%  &    816 & 11.3 &  159 \\
2023 &   326 &  13.9\,\% &  5{,}321 & 16.3 &  478 \\
2024 &   684 &  29.1\,\% & 11{,}891 & 17.4 &  636 \\
2025 &   885 &  37.6\,\% & 11{,}007 & 12.4 &  863 \\
2026 &   386 &  16.4\,\% &  4{,}872 & 12.6 &  552 \\
\midrule
\textbf{total} & \textbf{2{,}353} & \textbf{100\,\%} & \textbf{33{,}907} & \textbf{14.4} & \textbf{1{,}338} \\
\bottomrule
\end{tabular}
\caption{Per-year statistics for the full release. ``Unique kws''
counts distinct registry-supplied keyword strings observed in
documents published that year (any year-membership; rows do not sum).
The 2024 spike in average questions per document is driven by long
IP-Box filings (Tab.~\ref{tab:topic-distribution}).}
\label{tab:per-year}
\end{table}

Topical structure is summarised by the registry's keyword field
\textit{as published} --- no manual taxonomy or post-hoc clustering.
A document is counted toward a keyword whenever that exact string
appears in its keyword list. The top-30 keywords by document
frequency (Tab.~\ref{tab:topic-distribution}) cover $75.1\%$ of the
release; extending to top-50 covers $84.9\%$ and top-100 covers
$91.2\%$. The full vocabulary contains $1{,}338$ distinct keywords
with a long tail: median document carries 4 keywords, max~21.

\begin{table}[!htb]
\centering
\footnotesize
\begin{tabular}{@{}rrl@{}}
\toprule
\textbf{docs} & \textbf{\%} & \textbf{keyword} \\
\midrule
475 & 20.2\,\% & IP Box \\
295 & 12.5\,\% & stawka-stawki podatku-stawka preferencyjna podatku \\
280 & 11.9\,\% & oprogramowanie \\
242 & 10.3\,\% & komputery-program komputerowy \\
238 & 10.1\,\% & dzia\l{}alno\'s\'c-dzia\l{}alno\'s\'c gospodarcza \\
235 & 10.0\,\% & dzia\l{}alno\'s\'c-dzia\l{}alno\'s\'c badawczo-rozwojowa \\
210 &  8.9\,\% & nieruchomo\'sci-sprzeda\.z nieruchomo\'sci \\
174 &  7.4\,\% & czynno\'sci-czynno\'sci niepodlegaj\k{a}ce opodatkowaniu \\
163 &  6.9\,\% & nieruchomo\'sci-nieruchomo\'s\'c niezabudowana \\
147 &  6.2\,\% & odliczenia-prawo do odliczenia \\
144 &  6.1\,\% & podatnik-podatnik podatku od towar\'ow i us\l{}ug \\
127 &  5.4\,\% & zwolnienie-zwolnienie przedmiotowe \\
121 &  5.1\,\% & podatnik \\
111 &  4.7\,\% & przedsi\k{e}biorstwa-zorganizowana cz\k{e}\'s\'c przedsi\k{e}biorstwa \\
110 &  4.7\,\% & prace-prace badawczo-rozwojowe \\
105 &  4.5\,\% & koszt-koszty uzyskania przychod\'ow \\
 98 &  4.2\,\% & grunty-grunt niezabudowany \\
 94 &  4.0\,\% & zwolnienie \\
 92 &  3.9\,\% & koszt-koszty kwalifikowane \\
 89 &  3.8\,\% & nieruchomo\'sci-nieruchomo\'s\'c zabudowana \\
 81 &  3.4\,\% & opodatkowanie \\
 80 &  3.4\,\% & gmina \\
 78 &  3.3\,\% & czynno\'sci-czynno\'sci podlegaj\k{a}ce opodatkowaniu \\
 70 &  3.0\,\% & prawa-prawa autorskie \\
 68 &  2.9\,\% & maj\k{a}tek-maj\k{a}tek prywatny \\
 65 &  2.8\,\% & dzia\l{}ki \\
 63 &  2.7\,\% & rycza\l{}t-rycza\l{}t ewidencjonowany \\
 59 &  2.5\,\% & umowa-umowa przedwst\k{e}pna \\
 54 &  2.3\,\% & dzier\.zawa \\
 53 &  2.3\,\% & odliczenia-odliczenie podatku od towar\'ow i us\l{}ug \\
\bottomrule
\end{tabular}
\caption{Top-30 keywords by document frequency (registry strings
verbatim, no manual taxonomy or normalisation). The marginal column
\textbf{\%} is the fraction of the full release ($n=2{,}353$) that
carries that keyword; rows do not sum because a document typically
carries 4 keywords. The single most common keyword is \textit{IP Box}
($20.2\%$), which co-occurs strongly with \textit{stawka preferencyjna}
($12.5\%$), \textit{oprogramowanie} ($11.9\%$), \textit{program
komputerowy} ($10.3\%$) and \textit{dzia\l{}alno\'s\'c badawczo-rozwojowa}
($10.0\%$) --- together forming the IP-Box-relief lexicon. A second
cluster covers real-estate VAT (\textit{sprzeda\.z nieruchomo\'sci},
\textit{grunt niezabudowany}, \textit{nieruchomo\'s\'c zabudowana});
a third covers VAT-exempt activities (\textit{czynno\'sci niepodlegaj\k{a}ce
opodatkowaniu}, \textit{zwolnienie przedmiotowe}). The 2024 share of
\textit{IP Box} reaches $40\%$ of that year's documents (up from
$4\%$ in 2022 and falling to $10\%$ in 2025--2026) --- a
regulatory-cycle artefact rather than a sampling bias.}
\label{tab:topic-distribution}
\end{table}

% =====================================================================
% Evaluation procedure
% =====================================================================

\section{Evaluation Procedure (OPUS-GENERATED-TO-REFINE)}
\label{app:evaluation}

This appendix specifies the matching, scoring and leakage-control
machinery used in Tab.~\ref{tab:ablation-routes} (extraction) and in
the downstream question-recreation evaluation. Two design choices
deserve justification: a four-tier normalisation cascade for span
matching (rather than exact-string equality), and a temporally
constrained few-shot/RAG context for the generator.

\paragraph{Span matching: four-tier normalisation cascade.}
A predicted span $p$ matches a gold span $g$ when any tier of the
following cascade equates them; the cascade is calibrated against
five development docs and validated on twenty held-out docs.
\textbf{Tier 1} ($n$): collapse runs of whitespace to a single space,
fold Unicode quote/dash variants (\textit{„''"}, \textit{–—}) to
ASCII, lowercase. \textbf{Tier 2} ($n_{\text{strict}}$): strip
\emph{all} whitespace plus tier-1 punctuation/case folding --- this
salvages OCR artefacts such as the typo ``\textit{C zy}'' in
gold doc 641645, where the model emitted ``\textit{Czy}'' against a
broken-spacing source. \textbf{Tier 3}
($n_{\text{no\_prefix}}$): tier 2 plus removal of a leading numeric
or alphabetic list marker (\texttt{1)}, \texttt{2.}, \texttt{a)},
\texttt{Pytanie~1:}, \texttt{Ad.~3.}, \texttt{Na pytanie:});
calibrated against doc 599923 where the source carried ``\texttt{2)
Czy~\ldots?}'' on a separate line and the model dropped the
``\texttt{2)}'' marker. \textbf{Tier 4} ($n_{\text{no\_trail}}$):
tier 2 plus stripping trailing \texttt{?\,:\,.\,,\,;\,-\,–\,—};
calibrated against doc 553467 where the gold context ended
``\textit{wskazać:}'' and the model emitted ``\textit{wskazać}''.
A predicted text is accepted as text-matched against gold $g$ if
$n(p) = n(g)$, else $n_{\text{no\_prefix}}(p) = n_{\text{no\_prefix}}(g)$,
else $n_{\text{no\_trail}}(p) = n_{\text{no\_trail}}(g)$
(tier 2 alone is not used as the primary text-match channel because it
is too aggressive --- it would equate two questions that differ only
in spacing). Substring matching is deliberately excluded: it would
credit the model when it emits one of $N$ glued sub-questions or
truncates after the first ``?'', and those are real failures, not
scoring artefacts.

\paragraph{Context match.}
Once a text-match is established, the predicted context $p_c$ must
match the corresponding gold context $g_c$. Context applies the same
four-tier cascade as text. In addition, a single \emph{95\%
containment} rule is applied: if $|n(p_c)| \geq 30$ and
$|n(p_c)|/|n(g_c)| \geq 0.95$ (or vice versa) and the shorter is a
substring of the longer, the contexts are accepted as matched. This
forgives boundary truncation of a few words at either end (the model
sometimes drops a leading ``\textit{Wnioskodawca wyjaśnił, że:}'' from
a long preamble) without forgiving wholesale lead-in omission, which
would break the sibling-equality invariant from prompt rules R5/R6.

\paragraph{Multi-match handling.}
A single text can legitimately appear under two distinct contexts in
the source --- e.g.\ the noun-list \textit{organizacyjnej /
finansowej / funkcjonalnej} repeated twice under different parents
in doc 603682. The matcher therefore indexes gold by text and tries
\emph{every} gold entry sharing the predicted text in turn; the
prediction is strict-matched if its context lines up with \emph{any}
of those entries. Without this, the matcher would force a deterministic
pairing that the verbatim-substring contract does not require.

\paragraph{Aggregated metrics.}
Let $K$ be the post-filter kept count, $G$ the gold count, $T$
text-matches and $S$ strict matches.
$\text{Recall} = T/G$, $P_{\text{text}} = T/K$, and
$P_{\text{strict}} = S/K$. The third precision channel
$P_{\text{ctx}}$ is the conditional rate $S/T$: of predictions
whose text matched, what fraction also matched context. Splitting
out $P_{\text{ctx}}$ separates context drift (text right, parent
wrong --- typically R6 violations) from text drift (paraphrase or
hallucinated span --- R1 violations); the two failure modes call for
different fixes (rule wording vs.\ post-filter strictness) and
collapsing them into a single number obscures that signal. The
unweighted aggregate row in Tab.~\ref{tab:ablation-routes} is the
mean of the per-route per-document counts, not the per-route
arithmetic mean of percentages, so that one document's $K=0$ does
not contribute to the precision aggregate.

\paragraph{Few-shot leakage control.}
The unified six-document few-shot bank used by both extraction
prompts is excluded from every evaluation cohort; the ablation
sampler silently filters those document IDs from the candidate pool
and the dev-iteration harness fails with an explicit leakage warning
if a developer requests one. The same exclusion holds for the
downstream question-recreation evaluator. The bank itself was
selected to span the principal Polish tax instruments under which
KIS issues clarifications, so that no single instrument category
biases extraction behaviour: two VAT cases (one real-estate
transaction, doc 482229, 2022-02-28; one \textit{gmina}
local-government applicant, doc 484231, 2022-03-04), one
corporate-income-tax (CIT) financial-subsidy case (doc 601721,
2024-08-30), one personal-income-tax (PIT) international-residence
case (doc 634282, 2025-04-04), one IP-Box case (preferential 5\,\%
rate for software/IP income; doc 569200, 2023-11-08, which also
supplies the long-form scale signal at $\sim$26\,k characters), and
one general business-activity classification (\textit{działalność
gospodarcza}) case (doc 641645, 2025-05-28). The publication-date
span of the bank (2022-02 -- 2025-05) sits comfortably inside the
Verified split's 2021-12 -- 2026-04 window
(Tab.~\ref{tab:dataset-splits}), so excluding the bank does not
remove the temporal extremes of the evaluation cohort.

\paragraph{Temporal context for question generation.}
The downstream generator (\S~\ref{app:extraction:sampling}'s
companion task: producing clarifying questions from a factual state)
runs over a Polish-tax retrieval-augmented stack with two corpora:
historical KIS interpretations and statutory law articles. Both are
filtered by publication date before being shown to the model.
\textbf{Similar-doc retrieval}: the corpus is masked to
$\{d : \texttt{publication\_date}(d) < \texttt{publication\_date}(q)\}$
where $q$ is the query document; an interpretation that did not yet
exist when the office wrote the clarification cannot be retrieved
as evidence for it. \textbf{Statutory retrieval}: each article
carries an \texttt{effective\_from}/\texttt{effective\_to} pair and
is masked to articles in force on $\texttt{publication\_date}(q)$,
i.e.\ $\texttt{effective\_from} \leq d_q < \texttt{effective\_to}$.
This excludes both repealed-and-replaced articles (current at
evaluation time but not at $d_q$) and provisions that came into
force after $d_q$ --- the latter is the dominant leakage channel for
recent docs, since most law amendments since 2020 appear in the
training data of the LLMs we evaluate. \textbf{Few-shot examples}
for the generator are drawn from the gold pool and excluded if their
\texttt{doc\_id} appears in the evaluation cohort --- a separate,
identity-level guard from the publication-date filter, which catches
the case where a doc \emph{is} in the cohort but happened to be
published earlier than another cohort member. Together these three
filters make the retrieval context the model would have seen
\emph{at the time of writing} doc $q$, isolating model capability
from temporal leakage.

% =====================================================================
% Prompt listings
% =====================================================================

\newcommand{\rulepair}[2]{%
\par\noindent
\begin{minipage}[t]{0.485\linewidth}\raggedright #1\end{minipage}%
\hspace{1em}%
\begin{minipage}[t]{0.485\linewidth}\raggedright #2\end{minipage}%
\par\smallskip}

\begin{promptbox}{Extraction prompt --- short supplementary sections (small route)}
\label{fig:prompt-small}
\scriptsize
\linespread{0.95}\selectfont
\selectlanguage{polish}\sloppy\setlength{\emergencystretch}{6em}\hyphenpenalty=50\tolerance=9999\hbadness=10000

\rulepair{\textbf{Polish (deployed prompt)}}{\textbf{English (reading aid, not deployed)}}
\par\hrule height 0.4pt\par\smallskip

\rulepair{%
Ekstrahujesz pytania urzędu z sekcji uzupełnień polskich interpretacji podatkowych KIS (krótkie dokumenty, do~$\sim$10k znaków), podanej w postaci czystego tekstu.\par\smallskip
Wynik: lista \texttt{\{question, context\}}.\par
--- \texttt{question} --- dosłowny tekst pytania (bez markera listy \texttt{a.}/\texttt{b.}/\texttt{\textbullet})\par
--- \texttt{context} --- dosłowny tekst preambuły dla podpytań; \texttt{""} dla pytań nadrzędnych\par\smallskip
Sygnały pozycji listy w tekście:\par
--- linia zaczyna się od krótkiego markera listy: \texttt{a)}/\texttt{b)}/\texttt{a.}/\texttt{b.}/\texttt{i.}/\texttt{ii.}/\texttt{Ø}/\texttt{\textbullet}/\texttt{*}/\texttt{-}/\texttt{--}/\texttt{>} (jednoliterowy/jednosymbolowy znak na POCZĄTKU linii, nie wewnątrz zdania)\par
--- LUB linia z wcięciem (spacje/tab) bezpośrednio pod akapitem kończącym się \texttt{:}, \texttt{tj.:} lub \texttt{tzn.:}\par
--- LUB linia bezpośrednio po cytowanym fragmencie kończącym się \texttt{»:}, \texttt{":} lub \texttt{„\ldots'':} (cytowany fragment pełni rolę preambuły --- pozycje listy poniżej traktuj jako podpytania, a \texttt{context} = cały cytowany fragment)\par\smallskip
Marker listy~$\neq$~numeracja nadrzędne: \texttt{1)}, \texttt{2)}, \texttt{Pytanie~1:} to CZĘŚĆ TREŚCI numerowanego pytania na poziomie root (reguła~3), NIE marker do pominięcia. Pomijaj jako marker tylko krótkie symbole \texttt{a.}/\texttt{b.}/\texttt{\textbullet}/\texttt{-}/\texttt{--}/\texttt{>} itp.\par\smallskip
Reguły:%
}{%
You extract office questions from the supplementary sections of Polish KIS tax interpretations (short documents, up to~$\sim$10k characters), provided as plain text.\par\smallskip
Output: a list of \texttt{\{question, context\}}.\par
--- \texttt{question} --- the verbatim text of the question (without the list marker \texttt{a.}/\texttt{b.}/\texttt{\textbullet}).\par
--- \texttt{context} --- the verbatim preamble for sub-questions; \texttt{""} for top-level questions.\par\smallskip
List-position signals in the text:\par
--- a line begins with a short list marker: \texttt{a)}/\texttt{b)}/\texttt{a.}/\texttt{b.}/\texttt{i.}/\texttt{ii.}/\texttt{Ø}/\texttt{\textbullet}/\texttt{*}/\texttt{-}/\texttt{--}/\texttt{>} (a one-letter/one-symbol token at the START of a line, not inside a sentence)\par
--- OR an indented line (spaces/tab) directly beneath a paragraph ending in \texttt{:}, \texttt{tj.:} or \texttt{tzn.:}\par
--- OR a line directly after a quoted fragment ending in \texttt{»:}, \texttt{":} or \texttt{„\ldots'':} (the quoted fragment plays the role of the preamble --- treat the list items below as sub-questions, with \texttt{context} = the entire quoted fragment)\par\smallskip
List marker~$\neq$~top-level numbering: \texttt{1)}, \texttt{2)}, \texttt{Pytanie~1:} are PART OF the content of a numbered root-level question (rule~3), NOT a marker to drop. Drop as a marker only short symbols \texttt{a.}/\texttt{b.}/\texttt{\textbullet}/\texttt{-}/\texttt{--}/\texttt{>} and the like.\par\smallskip
Rules:%
}

\rulepair{%
1.~Tekst dosłowny: każdy \texttt{question} i \texttt{context} jest dosłownym fragmentem tekstu źródłowego, słowo w słowo. Bez parafrazy, bez uzupełniania.%
}{%
1.~Verbatim text: every \texttt{question} and \texttt{context} is a verbatim fragment of the source text, word for word. No paraphrase, no completion.%
}

\rulepair{%
2.~Pomijaj odpowiedzi wnioskodawcy (\texttt{Odp.:}/\texttt{Odpowiedź:}/„Stanowisko Wnioskodawcy''). Pytanie kończy się na ostatnim \texttt{?} lub \texttt{.} przed odpowiedzią lub kolejną numerowaną pozycją.%
}{%
2.~Skip applicant answers (\texttt{Odp.:}/\texttt{Odpowiedź:}/``Stanowisko Wnioskodawcy''). The question ends at the last \texttt{?} or \texttt{.} before the answer or the next numbered item.%
}

\rulepair{%
3.~Każda nowa numerowana pozycja (\texttt{1)}, \texttt{2)}, \texttt{Pytanie~1:}) wprowadza nowe pytanie. Nie łącz dwóch numerowanych pytań.\par\smallskip
Wrapper \texttt{Na pytanie:} / \texttt{Na zadane pytanie:} / \texttt{W~odpowiedzi na pytanie:} / \texttt{Pytanie nr~X brzmiało:} --- to znacznik przytoczonego pytania urzędu. Tekst od końca wrappera do najbliższego \texttt{Odpowiedział Pan:} / \texttt{Odpowiedział Pan;} / \texttt{Odp.:} / \texttt{Odpowiedź:} to TREŚĆ pytania urzędu --- wyemituj ją jako osobne NADRZĘDNE pytanie (z \texttt{context~=~""}). Jeśli ta treść zawiera kilka znaków \texttt{?}, stosuj regułę~9 (multi-\texttt{?} proza) wewnątrz niej. Jeśli kończy się preambułą \texttt{:} z listą bulletów, stosuj regułę~5 (podpytania z preambułą jako \texttt{context}). Sam wrapper nie pojawia się w wyniku.%
}{%
3.~Every new numbered item (\texttt{1)}, \texttt{2)}, \texttt{Pytanie~1:}) introduces a new question. Do not merge two numbered questions.\par\smallskip
The wrapper \texttt{Na pytanie:} / \texttt{Na zadane pytanie:} / \texttt{W~odpowiedzi na pytanie:} / \texttt{Pytanie nr~X brzmiało:} is a marker of a quoted office question. The text from the end of the wrapper up to the nearest \texttt{Odpowiedział Pan:} / \texttt{Odpowiedział Pan;} / \texttt{Odp.:} / \texttt{Odpowiedź:} is the CONTENT of the office question --- emit it as a separate TOP-LEVEL question (with \texttt{context~=~""}). If that content contains several \texttt{?} marks, apply rule~9 (multi-\texttt{?} prose) inside it. If it ends with a preamble \texttt{:} followed by a bulleted list, apply rule~5 (sub-questions with the preamble as \texttt{context}). The wrapper itself does not appear in the output.%
}

\rulepair{%
4.~Preambuła to akapit kończący się znakiem \texttt{:}, \texttt{tj.:} lub \texttt{tzn.:}, bezpośrednio poprzedzający pozycje listy. Akapit kończący się BEZ \texttt{:} ALE bezpośrednio poprzedzający pozycje listy z markerem (\texttt{-}, \texttt{\textbullet}, \texttt{a)}, \ldots) RÓWNIEŻ traktuj jako preambułę --- emituj jej treść jako \texttt{context} podpytań, NIE jako pytanie. Nie emituj preambuły jako pytania, nawet jeśli zaczyna się od \texttt{Czy}/\texttt{Co}/\texttt{Jak}.\par\smallskip
\textbf{NARRACJA $\neq$ preambuła}: zdania wprowadzające typu „W~uzupełnieniu wniosku Wnioskodawca udzielił odpowiedzi:'' --- to NARRACJA, nie preambuła. Pytania pod nimi to NADRZĘDNE pytania urzędu (każde ma \texttt{context~=~""}), NIE podpytania. Sygnał narracji: zdanie opisujące CO ZROBIŁ wnioskodawca/organ (czasownik \texttt{wskazali}, \texttt{udzielił}, \texttt{odpowiedzieli}), zakończone \texttt{:} ale NIE wprowadzające numerowanej/literowanej listy.\par\smallskip
Wyjątek (preambuła zawierająca pełne pytanie nadrzędne): jeśli preambuła ma postać \texttt{Czy~X?; jeżeli tak, proszę wskazać:}, to część przed \texttt{;} (zdanie pytajne) wyemituj jako pytanie nadrzędne z \texttt{context~=~""}, a cała preambuła służy jako \texttt{context} podpytań.\par\smallskip
Kontrola po wygenerowaniu: WYRZUĆ każdą pozycję wyniku której pole \texttt{question} (po stripie spacji) kończy się \texttt{:} --- to preambuły.%
}{%
4.~A preamble is a paragraph ending in \texttt{:}, \texttt{tj.:} or \texttt{tzn.:}, directly preceding list items. A paragraph ending WITHOUT \texttt{:} BUT directly preceding marked list items (\texttt{-}, \texttt{\textbullet}, \texttt{a)}, \ldots) ALSO counts as a preamble --- emit its text as the \texttt{context} of the sub-questions, NOT as a question. Do not emit a preamble as a question even when it begins with \texttt{Czy}/\texttt{Co}/\texttt{Jak}.\par\smallskip
\textbf{NARRATION $\neq$ preamble}: introductory sentences such as ``W~uzupełnieniu wniosku Wnioskodawca udzielił odpowiedzi:'' are NARRATION, not preambles. The questions beneath them are TOP-LEVEL office questions (each with \texttt{context~=~""}), NOT sub-questions. Narration signal: a sentence describing WHAT the applicant/office DID (verb \texttt{wskazali}, \texttt{udzielił}, \texttt{odpowiedzieli}), ending in \texttt{:} but NOT introducing a numbered/lettered list.\par\smallskip
Exception (preamble containing a full parent question): if the preamble has the form \texttt{Czy~X?; jeżeli tak, proszę wskazać:}, emit the part before \texttt{;} (the interrogative clause) as a parent question with \texttt{context~=~""}; the entire preamble serves as the \texttt{context} of the sub-questions.\par\smallskip
Post-generation check: DROP any output item whose \texttt{question} field (after stripping whitespace) ends in \texttt{:} --- those are preambles.%
}

\rulepair{%
5.~Pozycje listy bezpośrednio pod preambułą emituj jako oddzielne podpytania, z treścią preambuły jako \texttt{context}. Wyemituj każdą pozycję; marker listy jest pomijany w \texttt{question}.%
}{%
5.~Emit list items directly below a preamble as separate sub-questions, with the preamble's text as \texttt{context}. Emit every item; the list marker is dropped from \texttt{question}.%
}

\rulepair{%
6.~Podpytania pod tą samą preambułą mają identyczny \texttt{context} znak po znaku --- bez skracania, bez wybierania innego zdania. Żaden \texttt{context} nie może być ucięty w środku słowa.%
}{%
6.~Sub-questions under the same preamble share an identical \texttt{context}, character for character --- no truncation, no picking a different sentence. No \texttt{context} may be cut in the middle of a word.%
}

\rulepair{%
7.~Wyjątek od reguły~5 (wyliczenie rzeczownikowe): jeśli wszystkie pozycje pod preambułą to krótkie frazy rzeczownikowe (mniej niż~6 słów, brak czasownika osobowego), to całość jest jednym pytaniem nadrzędnym z wyliczeniem wewnątrz --- nie dziel.%
}{%
7.~Exception to rule~5 (noun enumeration): if all items below the preamble are short noun phrases (fewer than~6 words, no finite verb), then the whole block is a single parent question with the enumeration inside --- do not split.%
}

\rulepair{%
8.~Spójniki \texttt{lub}/\texttt{oraz}/\texttt{albo}/\texttt{i}/\texttt{czy~też}, lub przecinek-łącznik (\texttt{X, Y lub Z?}), łączące alternatywy rzeczownikowe w obrębie jednego zdania pytajnego, oznaczają jedno pytanie. Sygnały podziału (zachowaj podział): drugie zdanie ma własny podmiot+orzeczenie, lub zaczyna się czasownikiem osobowym po koordynatorze.%
}{%
8.~Connectors \texttt{lub}/\texttt{oraz}/\texttt{albo}/\texttt{i}/\texttt{czy~też}, or a linking comma (\texttt{X, Y lub Z?}), joining noun-phrase alternatives within one interrogative clause indicate a single question. Split signals (keep the split): the second clause has its own subject+predicate, or begins with a finite verb after the coordinator.%
}

\rulepair{%
9.~Wiele \texttt{?} w jednym akapicie (proza): \textbf{NIE DZIEL} gdy drugie zdanie pasuje do jednego z wzorców kontynuacji: warunek \texttt{Jeśli/Jeżeli/Gdy}\,+\,odniesienie do poprzedniej odpowiedzi (\texttt{tak}/\texttt{nie}) lub anafora bez własnego rzeczownika; imperatyw doprecyzowujący (\texttt{Proszę}/\texttt{Należy}); spójnik wyjaśniający (\texttt{tj.}/\texttt{tzn.}/\texttt{np.}/\texttt{w~szczególności}); wtrącenie w nawiasie. \textbf{DZIEL} gdy drugie zdanie wprowadza NOWY samodzielny temat (własny rzeczownik tematyczny + własny podmiot+orzeczenie). Wspólny adresat (\texttt{Pan}/\texttt{Pani}/\texttt{Wnioskodawca}) nie tworzy zależności. Reguła~5 ma pierwszeństwo nad regułą~9.%
}{%
9.~Multiple \texttt{?} in a single paragraph (prose): \textbf{DO NOT SPLIT} when the second clause matches a continuation pattern: a condition \texttt{Jeśli/Jeżeli/Gdy}\,+\,a back-reference to the previous answer (\texttt{tak}/\texttt{nie}) or an anaphor without its own noun; a clarifying imperative (\texttt{Proszę}/\texttt{Należy}); an explanatory connector (\texttt{tj.}/\texttt{tzn.}/\texttt{np.}/\texttt{w~szczególności}); a parenthetical insertion. \textbf{SPLIT} when the second clause introduces a NEW self-contained topic (its own topic noun + its own subject+predicate). A shared addressee (\texttt{Pan}/\texttt{Pani}/\texttt{Wnioskodawca}) does not create a dependency. Rule~5 takes precedence over rule~9.%
}

\rulepair{%
10.~\texttt{context} musi być pełnym kontekstem semantycznym; krótkie łączniki (\texttt{proszę wyjaśnić:}/\texttt{należy wskazać:}/\texttt{Zatem proszę wskazać:}) są niewystarczające. Gdy poprzedzający akapit pasuje do wzorca wprowadzającego (konstrukcje odniesienia \texttt{W~związku z\ldots}/\texttt{W~opisie sprawy\ldots}/\texttt{We wniosku\ldots}/\texttt{Odnosząc się do\ldots}; czasownik referencyjny \texttt{wskazali/podali/oświadczyli Państwo/Pan/Pani}; lub cytat treści wniosku) --- włącz CAŁY ten akapit (i poprzednie tej samej grupy wprowadzającej) do \texttt{context}. Reguła~10 ma pierwszeństwo nad regułą~4.%
}{%
10.~The \texttt{context} must be the full semantic context; short connectives (\texttt{proszę wyjaśnić:}/\texttt{należy wskazać:}/\texttt{Zatem proszę wskazać:}) are insufficient. When the preceding paragraph matches the lead-in pattern (reference constructions \texttt{W~związku z\ldots}/\texttt{W~opisie sprawy\ldots}/\texttt{We wniosku\ldots}/\texttt{Odnosząc się do\ldots}; a referential verb \texttt{wskazali/podali/oświadczyli Państwo/Pan/Pani}; or a quote of application content), include the WHOLE such paragraph (and any earlier paragraphs in the same lead-in group) in \texttt{context}. Rule~10 takes precedence over rule~4.%
}

\rulepair{%
11.~Pomiń \texttt{question} TYLKO gdy WSZYSTKO: brak \texttt{?}; brak słowa pytajnego/imperatywu (\texttt{Czy}/\texttt{Co}/\texttt{Jak}/\texttt{Jakie}/\texttt{Który}/\texttt{Kiedy}/\texttt{Gdzie}/\texttt{Kto}/\texttt{Ile}/\texttt{Proszę}/\texttt{Prosimy}/\texttt{Należy}/\texttt{Winien}/\texttt{Wskazanie} oraz odmiany); nie jest wezwaniem warunkowym (\texttt{Jeżeli/Jeśli X to czy/proszę/należy Y}, \texttt{Wobec X\ldots}, \texttt{W~przypadku X\ldots}); oraz nie kończy się \texttt{;}/\texttt{,}. Czysta narracja lub cytat z wniosku bez wezwania jest pomijany.%
}{%
11.~Skip a \texttt{question} ONLY when ALL of: no \texttt{?}; no interrogative/imperative opener (\texttt{Czy}/\texttt{Co}/\texttt{Jak}/\texttt{Jakie}/\texttt{Który}/\texttt{Kiedy}/\texttt{Gdzie}/\texttt{Kto}/\texttt{Ile}/\texttt{Proszę}/\texttt{Prosimy}/\texttt{Należy}/\texttt{Winien}/\texttt{Wskazanie} and inflected forms); not a conditional summons (\texttt{Jeżeli/Jeśli X to czy/proszę/należy Y}, \texttt{Wobec X\ldots}, \texttt{W~przypadku X\ldots}); and does not end in \texttt{;}/\texttt{,}. Bare narration or a quote from the application without a summons is dropped.%
}

\rulepair{%
12.~Pusta lista jest dopuszczalna wyłącznie gdy sekcja zawiera tylko akapity zaczynające się od \texttt{Odp.:}/\texttt{Odpowiedź:}/„Stanowisko Wnioskodawcy'', lub samo oświadczenie wnioskodawcy bez pytań organu.%
}{%
12.~An empty list is allowed only when the section contains nothing but paragraphs beginning with \texttt{Odp.:}/\texttt{Odpowiedź:}/``Stanowisko Wnioskodawcy'', or a stand-alone applicant statement without office questions.%
}
\end{promptbox}
\noindent\textit{The left column reproduces the deployed prompt verbatim; the right column is an English version translated with Claude Opus~4.7 as a reading aid (not deployed).}

\begin{promptbox}{Extraction prompt --- long supplementary sections (large route, HTML-aware)}
\label{fig:prompt-large-html}
\scriptsize
\linespread{0.95}\selectfont
\selectlanguage{polish}\sloppy\setlength{\emergencystretch}{6em}\hyphenpenalty=50\tolerance=9999\hbadness=10000
\rulepair{\textbf{Polish (deployed prompt)}}{\textbf{English (reading aid, not deployed)}}
\par\hrule height 0.4pt\par\smallskip

\rulepair{%
Ekstrahujesz pytania urzędu z sekcji uzupełnień polskich interpretacji podatkowych KIS. Wejście: HTML MS Word z zachowaną strukturą akapitów i list.\par\smallskip
Wynik: lista \texttt{\{question, context\}}.\par
--- \texttt{question} --- dosłowny wyrenderowany tekst pytania, bez tagów HTML i bez markera listy. Encje HTML rozwiń (\texttt{\&nbsp;}$\to$spacja, \texttt{\&amp;}$\to$\texttt{\&}, \texttt{\&ndash;}$\to$\texttt{--}, \texttt{\&mdash;}$\to$\texttt{---}, \texttt{\&rsquo;}$\to$\texttt{'}, \texttt{\&ldquo;}$\to$„, \texttt{\&rdquo;}$\to$'').\par
--- \texttt{context} --- dosłowny wyrenderowany tekst preambuły dla podpytań; \texttt{""} dla pytań nadrzędnych.\par\smallskip
Sygnały pozycji listy w HTML (DOWOLNY z poniższych wystarcza --- alternatywa, nie koniunkcja):\par
--- \textbf{KLASA}: \texttt{<p>} z klasą stylu wskazującą wcięcie/pozycję listy (\texttt{MsoListParagraph*}, \texttt{MsoNormalCxSp*}, lub równoważna)\par
--- \textbf{STYL}: atrybut \texttt{style} z \texttt{margin-left:}~$>$~0 lub \texttt{text-indent:}~$<$~0 (wcięcie wiszące)\par
--- \textbf{TREŚĆ}: pierwsze widoczne znaki akapitu to krótki marker rozdziału (\texttt{a.}/\texttt{b.}/\texttt{1.}/\texttt{i.}/\texttt{Ø}/\texttt{\textbullet}/\texttt{*}/\texttt{·}/\texttt{-}/\texttt{--}) LUB obecny \texttt{<span style="mso-list:Ignore">marker</span>}\par\smallskip
\textbf{Atypowe dokumenty}: gdy klasa to \texttt{MsoBodyText}/\texttt{MsoNormal} ALE akapit ma \texttt{margin-left:}~$>$~0 lub \texttt{text-indent:}~$<$~0 lub treść zaczyna się markerem --- traktuj jako pozycję listy.\par\smallskip
Granica listy: kolejne pozycje (\texttt{...CxSpFirst $\to$ Middle $\to$ Last}, lub seria z tym samym \texttt{margin-left}) tworzą jedną listę. Lista kończy się gdy następny \texttt{<p>} ma inną klasę / inny \texttt{margin-left}, lub jest preambułą poprzedzającą KOLEJNĄ listę.\par\smallskip
Numeracja pytań nadrzędnych (\texttt{1)}/\texttt{2)}/\texttt{Pytanie~1:}/\texttt{Ad.~1.}/\texttt{pytanie nr~X ww.\ wezwania:}) NIE jest markerem listy do pominięcia --- to część treści pytania nadrzędnego (reguła~3).\par\smallskip
Reguły:%
}{%
You extract office questions from supplementary sections of Polish KIS tax interpretations. Input: MS Word HTML with paragraph and list structure preserved.\par\smallskip
Output: list of \texttt{\{question, context\}}.\par
--- \texttt{question} --- the verbatim rendered question text, with no HTML tags and no list marker. Expand HTML entities (\texttt{\&nbsp;}$\to$space, \texttt{\&amp;}$\to$\texttt{\&}, \texttt{\&ndash;}$\to$\texttt{--}, \texttt{\&mdash;}$\to$\texttt{---}, \texttt{\&rsquo;}$\to$\texttt{'}, \texttt{\&ldquo;}$\to$„, \texttt{\&rdquo;}$\to$'').\par
--- \texttt{context} --- the verbatim rendered preamble for sub-questions; \texttt{""} for parent questions.\par\smallskip
List-position signals in HTML (ANY one below suffices --- disjunction, not conjunction):\par
--- \textbf{CLASS}: \texttt{<p>} with a style class indicating indentation/list position (\texttt{MsoListParagraph*}, \texttt{MsoNormalCxSp*}, or equivalent)\par
--- \textbf{STYLE}: a \texttt{style} attribute with \texttt{margin-left:}~$>$~0 or \texttt{text-indent:}~$<$~0 (hanging indent)\par
--- \textbf{CONTENT}: the first visible characters of the paragraph form a short item marker (\texttt{a.}/\texttt{b.}/\texttt{1.}/\texttt{i.}/\texttt{Ø}/\texttt{\textbullet}/\texttt{*}/\texttt{·}/\texttt{-}/\texttt{--}) OR a \texttt{<span style="mso-list:Ignore">marker</span>} is present\par\smallskip
\textbf{Atypical documents}: when the class is \texttt{MsoBodyText}/\texttt{MsoNormal} BUT the paragraph has \texttt{margin-left:}~$>$~0 or \texttt{text-indent:}~$<$~0 or the content starts with a marker --- treat it as a list item.\par\smallskip
List boundary: consecutive items (\texttt{...CxSpFirst $\to$ Middle $\to$ Last}, or a series sharing the same \texttt{margin-left}) form a single list. The list ends when the next \texttt{<p>} has a different class / different \texttt{margin-left}, or is a preamble preceding ANOTHER list.\par\smallskip
Parent-question numbering (\texttt{1)}/\texttt{2)}/\texttt{Pytanie~1:}/\texttt{Ad.~1.}/\texttt{pytanie nr~X ww.\ wezwania:}) is NOT a list marker to drop --- it is part of the parent question's content (rule~3).\par\smallskip
Rules:%
}

\rulepair{%
1.~Dosłowność: \texttt{question} i \texttt{context} są dosłownymi fragmentami wyrenderowanego HTML, słowo w słowo. Bez parafrazy.%
}{%
1.~Verbatim: \texttt{question} and \texttt{context} are word-for-word fragments of the rendered HTML. No paraphrase.%
}

\rulepair{%
2.~Pomijaj odpowiedzi wnioskodawcy (\texttt{Odp.:}/\texttt{Odpowiedź:}/„Stanowisko Wnioskodawcy'') i czystą narrację. Pytanie kończy się ostatnim \texttt{?} lub \texttt{.} przed odpowiedzią lub kolejną numerowaną pozycją.%
}{%
2.~Skip applicant answers (\texttt{Odp.:}/\texttt{Odpowiedź:}/``Stanowisko Wnioskodawcy'') and bare narration. The question ends at the last \texttt{?} or \texttt{.} before an answer or the next numbered item.%
}

\rulepair{%
3.~Każda nowa numerowana pozycja nadrzędna wprowadza nowe pytanie. Wrapper przytoczonego pytania (\texttt{Na pytanie:}/\texttt{Pytanie nr~X brzmiało:}) wprowadza pytanie nadrzędne (\texttt{context~=~""}); sam wrapper jest pomijany. Wewnątrz przytoczonej treści: kilka \texttt{?} $\to$ reguła~9; \texttt{:} z listą bulletów $\to$ reguła~5.%
}{%
3.~Every new top-level numbered item starts a new question. A quoted-question wrapper (\texttt{Na pytanie:}/\texttt{Pytanie nr~X brzmiało:}) introduces a top-level question (\texttt{context~=~""}); the wrapper is dropped. Inside the quoted content: several \texttt{?} $\to$ rule~9; \texttt{:} followed by a bullet list $\to$ rule~5.%
}

\rulepair{%
4.~Preambuła = akapit kończący się \texttt{:}/\texttt{tj.:}/\texttt{tzn.:} lub niedokończonym zdaniem nad listą z markerem. Treść preambuły emituj jako \texttt{context} podpytań --- NIGDY jako osobne \texttt{question}. \textbf{Narracja $\neq$ preambuła}: zdania opisujące CZYNNOŚĆ wnioskodawcy/organu (\texttt{wskazali}/\texttt{udzielił}/\texttt{odpowiedzieli}) zakończone \texttt{:} przed PEŁNYMI numerowanymi pytaniami --- pomijaj.\par\smallskip
Wyjątek (preambuła z osadzonym pytaniem): \texttt{Czy~X?; jeżeli tak, proszę wskazać:} --- część przed \texttt{;}/\texttt{?}/\texttt{--} wyemituj jako pytanie nadrzędne z \texttt{context~=~""}; cała preambuła jest \texttt{context} podpytań pod nią.%
}{%
4.~Preamble = a paragraph ending in \texttt{:}/\texttt{tj.:}/\texttt{tzn.:} or an unfinished sentence above a marker-led list. Emit the preamble's text as the \texttt{context} of sub-questions --- NEVER as a stand-alone \texttt{question}. \textbf{Narration $\neq$ preamble}: sentences describing an ACTION by the applicant/the office (\texttt{wskazali}/\texttt{udzielił}/\texttt{odpowiedzieli}) ending with \texttt{:} before FULL numbered questions are NARRATION --- skip them.\par\smallskip
Exception (preamble with an embedded question): \texttt{Czy~X?; jeżeli tak, proszę wskazać:} --- emit the part before \texttt{;}/\texttt{?}/\texttt{--} as a top-level question with \texttt{context~=~""}; the whole preamble becomes the \texttt{context} of the sub-questions beneath it.%
}

\rulepair{%
5.~Pozycje listy pod preambułą emituj jako oddzielne podpytania z preambułą jako \texttt{context}. Marker listy pomijasz; \texttt{context} zawiera CAŁĄ wyrenderowaną treść preambuły, włącznie z numerowanym prefiksem (\texttt{Pytanie~1:}/\texttt{Ad.~1.}). Pytania nadrzędne mają \texttt{context~=~""}. Numeracja \texttt{1)}/\texttt{2)}/\texttt{3)} zamyka poprzednią listę: po liście pod \texttt{5)} \texttt{<p>6) Czy Y?</p>} jest KOLEJNYM pytaniem nadrzędnym. Dwie preambuły obok siebie obsługuj niezależnie.%
}{%
5.~Emit list items under a preamble as separate sub-questions with the preamble as their \texttt{context}. Drop the list marker; \texttt{context} contains the ENTIRE rendered preamble, including any numbered prefix (\texttt{Pytanie~1:}/\texttt{Ad.~1.}). Parent questions have \texttt{context~=~""}. Numbering \texttt{1)}/\texttt{2)}/\texttt{3)} closes the previous list: after the list under \texttt{5)}, \texttt{<p>6) Czy Y?</p>} is the NEXT parent question. Handle two adjacent preambles independently.%
}

\rulepair{%
6.~Podpytania jednej preambuły mają identyczny \texttt{context} znak po znaku. Nie ucinaj \texttt{context} w środku zdania ani słowa.%
}{%
6.~Sub-questions of one preamble share an identical \texttt{context} character-for-character. Do not truncate \texttt{context} mid-sentence or mid-word.%
}

\rulepair{%
7.~Wyjątek od reguły~5: gdy preambuła pyta o KOLEKCJĘ jako całość (\texttt{Czy w skład wchodzą:}/\texttt{Co obejmuje:}/yes-no o kolekcji) i wszystkie pozycje to krótkie frazy rzeczownikowe bez czasownika osobowego --- nie dziel; cały blok = JEDNO pytanie nadrzędne z wyliczeniem wewnątrz. Reguła~5 wraca gdy preambuła wymaga osobnej odpowiedzi per pozycja, jest negowana, lub którakolwiek pozycja jest pełnym zdaniem/kończy się \texttt{?}.%
}{%
7.~Exception to rule~5: when the preamble asks about a COLLECTION as a whole (\texttt{Czy w skład wchodzą:}/\texttt{Co obejmuje:}/yes-no about a collection) and all items are short noun phrases without a finite verb --- do not split; the whole block = ONE parent question with the enumeration inside. Rule~5 returns when the preamble demands a per-item answer, is negated, or any item is a full sentence / ends with \texttt{?}.%
}

\rulepair{%
8.~Spójniki alternatywy rzeczownikowej (\texttt{lub}/\texttt{oraz}/\texttt{albo}/\texttt{i}/\texttt{czy też}) lub przecinek-łącznik wewnątrz jednego zdania pytajnego = JEDNO pytanie. Sygnały podziału: druga część ma własny podmiot+orzeczenie, czasownik osobowy po koordynatorze, lub spójnik bezpośrednio przed \texttt{czy}/wh-słowem.%
}{%
8.~Noun-phrase coordinating conjunctions (\texttt{lub}/\texttt{oraz}/\texttt{albo}/\texttt{i}/\texttt{czy też}) or a comma-as-conjunction inside one interrogative sentence = ONE question. Split signals: the second clause has its own subject+predicate, a finite verb after the coordinator, or a conjunction directly before \texttt{czy}/a wh-word.%
}

\rulepair{%
9.~Wiele \texttt{?} w jednym \texttt{<p>}: NIE DZIEL gdy kontynuacja należy do tego samego pytania (warunek \texttt{Jeśli/Jeżeli/Gdy}\,+\,odniesienie/anafora bez własnego rzeczownika; imperatyw doprecyzowujący \texttt{Proszę}/\texttt{Należy}; spójnik wyjaśniający \texttt{tj.}/\texttt{tzn.}/\texttt{np.}/\texttt{w~szczególności}; wtrącenie w nawiasie). DZIEL gdy drugie zdanie wprowadza nowy rzeczownik tematyczny + własny podmiot+orzeczenie. Wspólny adresat (\texttt{Pan}/\texttt{Pani}/\texttt{Wnioskodawca}) nie tworzy zależności. Pierwszeństwo imperatywu doprecyzowującego nad analizą rzeczowników. Reguła~5 ma pierwszeństwo nad regułą~9.%
}{%
9.~Multiple \texttt{?} in one \texttt{<p>}: DO NOT SPLIT when the continuation belongs to the same question (condition \texttt{Jeśli/Jeżeli/Gdy}\,+\,reference/anaphor without its own noun; clarifying imperative \texttt{Proszę}/\texttt{Należy}; explanatory connector \texttt{tj.}/\texttt{tzn.}/\texttt{np.}/\texttt{w~szczególności}; parenthetical insertion). SPLIT when the second clause introduces a new topic noun + its own subject+predicate. A shared addressee (\texttt{Pan}/\texttt{Pani}/\texttt{Wnioskodawca}) does not create a dependency. The clarifying imperative takes precedence over the noun analysis. Rule~5 takes precedence over rule~9.%
}

\rulepair{%
10.~\texttt{context} musi być pełnym kontekstem semantycznym; krótkie łączniki (\texttt{proszę wyjaśnić:}/\texttt{należy wskazać:}/\texttt{Zatem proszę wskazać:}) są niewystarczające. Gdy poprzedzający akapit pasuje do wzorca wprowadzającego (przyimkowa konstrukcja odniesienia \texttt{W~związku z\ldots}/\texttt{W~opisie sprawy\ldots}/\texttt{We wniosku\ldots}/\texttt{Odnosząc się do\ldots}; czasownik mowy/wskazania w 2/3 osobie z podmiotem \texttt{Państwo}/\texttt{Pan}/\texttt{Pani}; cytat treści wniosku) --- włącz CAŁY ten akapit (i poprzednie tej samej grupy) do \texttt{context}. Reguła~10 $>$ reguła~4.%
}{%
10.~The \texttt{context} must be the full semantic context; short connectives (\texttt{proszę wyjaśnić:}/\texttt{należy wskazać:}/\texttt{Zatem proszę wskazać:}) are insufficient. When the preceding paragraph matches the lead-in pattern (a prepositional reference \texttt{W~związku z\ldots}/\texttt{W~opisie sprawy\ldots}/\texttt{We wniosku\ldots}/\texttt{Odnosząc się do\ldots}; a verb of speaking/indicating in 2nd/3rd person with subject \texttt{Państwo}/\texttt{Pan}/\texttt{Pani}; a quote of application content) --- include the WHOLE such paragraph (and any earlier paragraphs in the same group) in \texttt{context}. Rule~10 $>$ rule~4.%
}

\rulepair{%
11.~Pomiń \texttt{question} TYLKO gdy WSZYSTKO: brak \texttt{?}; brak słowa pytajnego/imperatywu (\texttt{Czy}/\texttt{Co}/\texttt{Jak}/\texttt{Jakie}/\texttt{Który}/\texttt{Kto}/\texttt{Gdzie}/\texttt{Kiedy}/\texttt{Ile}/\texttt{Proszę}/\texttt{Prosimy}/\texttt{Należy}/\texttt{Winien}/\texttt{Wskazanie} oraz odmiany; \texttt{czy} małą literą na początku zdania ZACHOWUJE status pytania); nie jest wezwaniem warunkowym z osadzonym pytaniem; nie kończy się \texttt{;}/\texttt{,}.%
}{%
11.~Skip a \texttt{question} ONLY when ALL of: no \texttt{?}; no interrogative/imperative opener (\texttt{Czy}/\texttt{Co}/\texttt{Jak}/\texttt{Jakie}/\texttt{Który}/\texttt{Kto}/\texttt{Gdzie}/\texttt{Kiedy}/\texttt{Ile}/\texttt{Proszę}/\texttt{Prosimy}/\texttt{Należy}/\texttt{Winien}/\texttt{Wskazanie} and inflected forms; lower-case \texttt{czy} at the start of a sentence STILL counts as a question); not a conditional summons with an embedded question; not ending in \texttt{;}/\texttt{,}.%
}

\rulepair{%
12.~Pusta lista jest dopuszczalna wyłącznie gdy sekcja zawiera same odpowiedzi wnioskodawcy lub samo oświadczenie bez pytań organu.%
}{%
12.~An empty list is allowed only when the section contains nothing but applicant answers or a stand-alone applicant statement without office questions.%
}
\end{promptbox}
\noindent\textit{The text-only large variant (\texttt{extractor\_supp\_large.yaml}) is a fallback when the cached HTML slice is shorter than the cleaned text; it drops the HTML-marker block and is otherwise identical.}
\noindent\textit{A text-only variant of this prompt is used as a
fallback when the source HTML is truncated; it drops the HTML-marker
block above and is otherwise identical. The left column reproduces
the deployed prompt verbatim; the right column is an English version
translated with Claude Opus~4.7 as a reading aid (not deployed).}

\begin{promptbox}{Zero-shot generator prompt --- general supplementary-question generator (no retrieved context)}
\label{fig:prompt-zero-shot-generator}
\scriptsize
\linespread{0.95}\selectfont
\selectlanguage{polish}\sloppy\setlength{\emergencystretch}{6em}\hyphenpenalty=50\tolerance=9999\hbadness=10000

\rulepair{\textbf{Polish (deployed prompt)}}{\textbf{English (reading aid, not deployed)}}
\par\hrule height 0.4pt\par\smallskip

\rulepair{%
Jesteś urzędnikiem organu interpretującego wzywającym wnioskodawcę do uzupełnienia opisu w sprawie.\par\smallskip
Twoim zadaniem jest wygenerowanie pytań uzupełniających (wezwania do uzupełnienia wniosku) na podstawie opisu stanu faktycznego / zdarzenia przyszłego przedstawionego przez podatnika. Celem pytań jest uzyskanie informacji niezbędnych do prawidłowej oceny skutków podatkowych opisanej sytuacji.\par\smallskip
Pytania mają wskazywać luki w opisie stanu faktycznego --- okoliczności istotne dla zastosowania powołanych przepisów.%
}{%
You are an officer of the interpreting authority calling on the applicant to supplement the description of the case.\par\smallskip
Your task is to generate supplementary questions (a call to supplement the application) based on the description of the factual state / future event presented by the taxpayer. The goal of the questions is to obtain information needed for a correct assessment of the tax consequences of the described situation.\par\smallskip
The questions should point to gaps in the factual-state description --- circumstances material to the application of the cited provisions.%
}

\rulepair{\textbf{=== STRUKTURA WIADOMOŚCI UŻYTKOWNIKA ===}}{\textbf{=== USER MESSAGE STRUCTURE ===}}

\rulepair{%
W treści zapytania użytkownika otrzymujesz:\par
--- ,,Stan faktyczny:'' --- opis stanu faktycznego / zdarzenia przyszłego przedstawiony przez podatnika.\par
--- ,,Pytania podatnika:'' (jeśli sekcja występuje) --- oryginalne pytania interpretacyjne, na które podatnik prosi o odpowiedź organu (sekcja ,,Pytania'' wniosku ORD-IN). Określają one ZAKRES sprawy: o jakie skutki podatkowe podatnik pyta i które normy prawne mają zostać zastosowane.%
}{%
The user message contains:\par
--- ``Stan faktyczny:'' (Factual State) --- the factual-state / future-event description presented by the taxpayer.\par
--- ``Pytania podatnika:'' (Taxpayer Questions, if the section is present) --- the original interpretive questions for which the taxpayer requests an answer (the ``Pytania'' section of the ORD-IN application). They define the SCOPE of the case: which tax consequences the taxpayer is asking about and which legal norms are to be applied.%
}

\rulepair{%
\textbf{Jak wykorzystywać ,,Pytania podatnika'':}\par
--- Traktuj je jako KLUCZ do priorytetyzacji pytań uzupełniających --- generuj przede wszystkim pytania o fakty potrzebne do udzielenia odpowiedzi na te pytania interpretacyjne.\par
--- Z pytań podatnika wynikają konkretne przepisy i instytucje prawne. Każda taka norma prawna ma własne przesłanki --- dla każdej sprawdź, czy stan faktyczny opisuje wszystkie elementy potrzebne do jej zastosowania, i pytaj o brakujące.\par
--- Jeśli pytanie podatnika powołuje konkretny artykuł, ustęp, punkt --- możesz cytować ten przepis w pytaniu uzupełniającym, gdy chcesz doprecyzować przesłankę.\par
--- NIE odpowiadaj na pytania podatnika ani nie oceniaj ich prawidłowości. Twoja rola to wyłącznie wyjaśnienie FAKTÓW potrzebnych do oceny.\par
--- Jeśli pytanie podatnika dotyczy zakresu, którego stan faktyczny nie obejmuje (np.\ pyta o jeden podatek, a opisuje sytuację tylko w zakresie innego) --- zadaj pytanie uzupełniające o brakujący wątek faktyczny.\par
--- Jeśli sekcja ,,Pytania podatnika'' jest pusta lub jej brak --- generuj pytania wyłącznie na podstawie stanu faktycznego, pokrywając wszystkie kierunki weryfikacji adekwatne do treści wniosku.%
}{%
\textbf{How to use ``Taxpayer Questions'':}\par
--- Treat them as the KEY to prioritising supplementary questions --- generate primarily questions about facts needed to answer those interpretive questions.\par
--- Specific provisions and legal institutions follow from the taxpayer's questions. Each such legal norm has its own preconditions --- for each, check whether the factual state describes all elements needed for its application, and ask about the missing ones.\par
--- If the taxpayer's question cites a specific article, section, or point --- you may quote that provision in the supplementary question when you want to clarify a precondition.\par
--- DO NOT answer the taxpayer's questions or judge their correctness. Your role is solely to clarify the FACTS needed for the assessment.\par
--- If the taxpayer's question concerns a scope not covered by the factual state (e.g.\ asks about one tax while the situation describes only another) --- ask a supplementary question about the missing factual thread.\par
--- If the ``Taxpayer Questions'' section is empty or missing --- generate questions based solely on the factual state, covering all verification directions adequate to the application's content.%
}

\rulepair{%
\textbf{Pytania powinny:}\par
--- Dotyczyć brakujących, niejasnych, niespójnych lub nadmiernie ogólnikowych informacji istotnych dla kwalifikacji prawnopodatkowej.\par
--- Być precyzyjne, konkretne i odwoływać się do treści wniosku (cytować lub parafrazować fragmenty, gdy to pomoże odbiorcy).\par
--- Dotyczyć wyłącznie kwestii faktycznych.\par
--- Dopuszczać jednoznaczną odpowiedź faktyczną; nie zapraszać do oceny subiektywnej ani opinii.\par
--- Być sformułowane w języku polskim, w stylu urzędowym (formalnym).\par
--- Zawierać pytania główne oraz podpytania (gdy pytanie wymaga doprecyzowania w wielu aspektach); pole \texttt{question} zawiera samą treść pytania, bez prefiksów typu ,,1)'', ,,a)'', ,,Pytanie:''. Numeracja wynika z kolejności w liście.%
}{%
\textbf{Questions should:}\par
--- Address missing, unclear, inconsistent, or overly generic information material to the tax-law qualification.\par
--- Be precise, concrete, and refer to the application's content (quote or paraphrase fragments where it helps the reader).\par
--- Address factual matters only.\par
--- Admit an unambiguous factual answer; not invite subjective judgement or opinion.\par
--- Be phrased in Polish, in an official (formal) register.\par
--- Contain parent questions and sub-questions (when a question needs clarification in multiple aspects); the \texttt{question} field carries only the question text, with no prefixes such as ``1)'', ``a)'', ``Pytanie:''. Numbering follows from list order.%
}

\rulepair{%
Format odpowiedzi: obiekt \texttt{\{"questions": [\ldots]\}} --- lista obiektów z polami:\par
--- \texttt{question}: treść pytania uzupełniającego.\par
--- \texttt{context}: dla pytania głównego --- pusty łańcuch \texttt{""}; dla podpytania --- dosłowna treść pytania głównego (znak po znaku identyczna we wszystkich podpytaniach jednej grupy).%
}{%
Response format: object \texttt{\{"questions": [\ldots]\}} --- a list of objects with fields:\par
--- \texttt{question}: the supplementary-question text.\par
--- \texttt{context}: for a parent question --- the empty string \texttt{""}; for a sub-question --- the verbatim text of the parent question (character-for-character identical across all sub-questions of one group).%
}

\rulepair{\textbf{=== WYBRANE KIERUNKI WERYFIKACJI STOSOWANE PRZEZ ORGAN PODATKOWY ===}}{\textbf{=== SELECTED VERIFICATION DIRECTIONS USED BY THE TAX AUTHORITY ===}}

\rulepair{%
Poniższe kierunki to przykładowy katalog typowych luk faktycznych --- nie zamknięta lista. Pytaj o każdą dostrzeżoną lukę, także taką, która nie pasuje do żadnego z kierunków. Nie powtarzaj informacji już wyczerpująco opisanych we wniosku.%
}{%
The directions below are an exemplary catalogue of typical factual gaps --- not a closed list. Ask about every gap you notice, including ones that do not fit any of the directions. Do not repeat information already exhaustively described in the application.%
}

\rulepair{%
1.~\textbf{PODMIOT I STATUS PODATNIKA} --- kto składa wniosek i w jakim charakterze. M.in.:\par
--- forma prawna działalności\par
--- rezydencja podatkowa, forma opodatkowania\par
--- status szczególny istotny dla sprawy\par
--- zakres i przedmiot prowadzonej działalności (kody PKD, faktyczny charakter działań)%
}{%
1.~\textbf{ENTITY AND TAXPAYER STATUS} --- who is filing the application and in what capacity. Among others:\par
--- legal form of activity\par
--- tax residency, form of taxation\par
--- special status material to the case\par
--- scope and subject of the activity (PKD codes, the actual nature of operations)%
}

\rulepair{%
2.~\textbf{OKRES I CHARAKTER WNIOSKU} --- czego dotyczy interpretacja w wymiarze czasowym. M.in.:\par
--- czy wniosek dotyczy stanu faktycznego (zdarzeń już zaistniałych) czy zdarzenia przyszłego (planowanego)\par
--- jakich lat podatkowych / okresów rozliczeniowych dotyczy\par
--- czy opisana sytuacja jest jednorazowa czy powtarzalna/ciągła\par
--- czy w sprawie toczy się lub toczyło postępowanie podatkowe, kontrola, lub wydano inną interpretację%
}{%
2.~\textbf{PERIOD AND NATURE OF THE APPLICATION} --- the temporal scope of the interpretation. Among others:\par
--- whether the application concerns a factual state (events already occurred) or a future event (planned)\par
--- which tax years / settlement periods it covers\par
--- whether the described situation is one-off or recurring/continuous\par
--- whether tax proceedings, an audit, or another interpretation is or was pending in the case%
}

\rulepair{%
3.~\textbf{STAN FAKTYCZNY --- SZCZEGÓŁY TRANSAKCJI / CZYNNOŚCI} --- dokładny opis tego, co się wydarzyło lub wydarzy. M.in.:\par
--- konkretne daty, kwoty, terminy, harmonogram zdarzeń\par
--- przedmiot transakcji --- cechy, klasyfikacja, przeznaczenie\par
--- sposób realizacji\par
--- motywacja ekonomiczna / cel biznesowy\par
--- źródło finansowania, sposób zapłaty, waluta rozliczenia%
}{%
3.~\textbf{FACTUAL STATE --- TRANSACTION / OPERATION DETAILS} --- a precise description of what happened or will happen. Among others:\par
--- specific dates, amounts, deadlines, schedule of events\par
--- subject of the transaction --- features, classification, intended use\par
--- manner of execution\par
--- economic motivation / business purpose\par
--- source of financing, method of payment, settlement currency%
}

\rulepair{%
4.~\textbf{STRONY I RELACJE} --- kto uczestniczy i w jakim powiązaniu. M.in.:\par
--- identyfikacja wszystkich stron\par
--- czy strony są podmiotami powiązanymi\par
--- tytuł prawny relacji między stronami (umowa cywilnoprawna, stosunek pracy, powiązanie korporacyjne)\par
--- zakres odpowiedzialności, ryzyka i praw poszczególnych stron%
}{%
4.~\textbf{PARTIES AND RELATIONSHIPS} --- who participates and in what relation. Among others:\par
--- identification of all parties\par
--- whether the parties are related entities\par
--- legal title of the relationship between the parties (civil-law contract, employment relationship, corporate affiliation)\par
--- scope of liability, risk, and rights of each party%
}

\rulepair{%
5.~\textbf{DOKUMENTACJA I EWIDENCJA} --- jak podatnik udokumentował stan faktyczny. M.in.:\par
--- umowy, aneksy, porozumienia --- ich treść w zakresie istotnym dla sprawy\par
--- dokumenty potwierdzające transakcję --- kto wystawia, komu, na jakich warunkach\par
--- ewidencje wymagane przepisami\par
--- oświadczenia, pełnomocnictwa, inne dokumenty potwierdzające istotne okoliczności%
}{%
5.~\textbf{DOCUMENTATION AND RECORDS} --- how the taxpayer documented the factual state. Among others:\par
--- contracts, annexes, agreements --- their content as relevant to the case\par
--- documents confirming the transaction --- who issues them, to whom, on what terms\par
--- records required by regulations\par
--- declarations, powers of attorney, other documents confirming material circumstances%
}

\rulepair{%
6.~\textbf{KWALIFIKACJA PRAWNA} --- elementy faktyczne potrzebne do zastosowania konkretnej normy. M.in.:\par
--- jeśli wniosek powołuje konkretny przepis (artykuł, ustęp, punkt) --- czy przedstawiony stan faktyczny wyczerpuje przesłanki tej normy\par
--- czy zostały spełnione warunki formalne (zgłoszenia, rejestracje, wpisy, decyzje)\par
--- czy nie zachodzą okoliczności wyłączające zastosowanie normy (wyjątki, limity, terminy prekluzyjne)\par
--- czy klasyfikacja proponowana przez podatnika znajduje oparcie w opisanym stanie faktycznym%
}{%
6.~\textbf{LEGAL QUALIFICATION} --- factual elements needed to apply a specific norm. Among others:\par
--- if the application cites a specific provision (article, section, point) --- whether the presented factual state satisfies the preconditions of that norm\par
--- whether the formal conditions have been met (notifications, registrations, entries, decisions)\par
--- whether no circumstances exclude the norm's application (exceptions, limits, preclusive deadlines)\par
--- whether the classification proposed by the taxpayer is supported by the described factual state%
}

\rulepair{%
7.~\textbf{SKUTKI I WZAJEMNE POWIĄZANIA} --- szersze następstwa czynności. M.in.:\par
--- jak opisana czynność wpływa na inne rozliczenia podatnika\par
--- czy opisana czynność nie stanowi elementu większej całości, której nie opisano we wniosku\par
--- czy podatnik korzystał / zamierza korzystać z ulg, zwolnień, odliczeń związanych z tą samą czynnością%
}{%
7.~\textbf{CONSEQUENCES AND INTERCONNECTIONS} --- broader downstream effects of the operation. Among others:\par
--- how the described operation affects the taxpayer's other settlements\par
--- whether the described operation is not part of a larger whole not described in the application\par
--- whether the taxpayer used or intends to use any reliefs, exemptions, or deductions tied to the same operation%
}

\rulepair{\textbf{=== ZASADY REDAKCYJNE ===}}{\textbf{=== EDITORIAL RULES ===}}

\rulepair{%
--- NIE zadawaj pytań prawnych (,,czy mam prawo do\ldots'', ,,czy moje stanowisko jest prawidłowe''). Pytaj o FAKTY.\par
--- NIE duplikuj informacji, które są już wyczerpująco opisane we wniosku. Każde pytanie musi ujawniać brakującą lub niejasną informację.\par
--- NIE zadawaj pytań niezwiązanych z meritum interpretacji. Każde pytanie powinno mieć znaczenie dla oceny skutków podatkowych.\par
--- ZACHOWUJ wszystkie istotne kwalifikatory faktyczne i prawne występujące we wniosku.\par
--- PRZY PYTANIACH ZŁOŻONYCH każde podpytanie to osobny rekord listy; pole \texttt{context} tego rekordu zawiera dosłowną treść pytania głównego.\par
--- PYTAJ O JEDNĄ RZECZ NA RAZ w podpytaniu. Jeśli temat wymaga wielu aspektów, rozbij na osobne podpytania (osobne rekordy listy z tym samym \texttt{context}).\par
--- JĘZYK URZĘDOWY: ,,Proszę wskazać\ldots'', ,,Proszę wyjaśnić\ldots'', ,,Czy\ldots'', ,,W jakim okresie\ldots'', ,,Na jakiej podstawie\ldots''.%
}{%
--- DO NOT ask legal questions (``do I have a right to\ldots'', ``is my position correct''). Ask about FACTS.\par
--- DO NOT duplicate information already exhaustively described in the application. Every question must surface a missing or unclear piece of information.\par
--- DO NOT ask questions unrelated to the merits of the interpretation. Every question must matter for the assessment of tax consequences.\par
--- PRESERVE all material factual and legal qualifiers present in the application.\par
--- FOR COMPOUND QUESTIONS, each sub-question is a separate list record; that record's \texttt{context} field carries the verbatim text of the parent question.\par
--- ASK ABOUT ONE THING AT A TIME in a sub-question. When the topic needs multiple aspects, split into separate sub-questions (separate list records sharing the same \texttt{context}).\par
--- OFFICIAL REGISTER: ``Proszę wskazać\ldots'' (Please indicate\ldots), ``Proszę wyjaśnić\ldots'' (Please explain\ldots), ``Czy\ldots'' (Does/Is\ldots), ``W jakim okresie\ldots'' (In what period\ldots), ``Na jakiej podstawie\ldots'' (On what basis\ldots).%
}

\end{promptbox}
\noindent\textit{The left column reproduces the deployed Polish prompt verbatim; the right column is an English version translated with Claude Opus~4.7 as a reading aid (not deployed). Source: \texttt{generator\_general.yaml}.}
